 % !TeX root = book.tex
\chapter{Prolog agents}\label{sec:prolog-agents}

\begin{quotation}
\noindent Imagine a world of \textbf{Prolog agents}, some useful, others playful, bringing joy to games and virtual worlds; some short-lived, others long-lived, some simple, others complex -- but maybe built from simpler ones. Written in Web Prolog, talking Web Prolog with other agents, using Web Prolog knowledge bases to guide their actions and conversations, making sure important capabilities of clever conversational agents, such as natural language understanding, knowledge representation, reasoning and real-time interaction, are accounted for.

\flushright\textit{Prolog agents -- the elevator pitch}
\end{quotation}

\vspace{1cm}

\noindent Chapter~\ref{sec:intro} introduced the notion of a Prolog agent, and Chapter~\ref{sec:web-prolog} the more precise concept of a Prolog actor -- the most elementary form of Prolog agent in the Prolog Trinity ecosystem. In this chapter, the concept of a Prolog agent is further developed and two important kinds of agents, Prolog toplevels and Prolog nodes, will be introduced and their roles in the ecosystem described.

The reason for referring to both actors and nodes as agents is to focus on their similarities. The similarity that makes us refer to both of them as Prolog agents is that they are both processes that are capable of talking Prolog to other agents. Also, they both live on the Prolog Web. There are of course differences as well. A node is typically long-lived. It is capable of serving many clients at the same time. It can be programmed, but only by its owner.

\section{Prolog agents in a nutshell}\label{sec:prolog-agents-nutshell}

In this book, an \emph{agent} is understood in a deliberately modest sense: a persistent software process that participates in ongoing interaction with an external environment.\footnote{In this spirit we follow \citet{russell2016artificial} in treating ``agent'' as an analytical lens rather than an absolute classification.}

\medskip
\noindent\textbf{Prolog agents.}
A \emph{Prolog agent} is a process on the Prolog Web whose interactions are expressed as Prolog terms: it receives goals or messages and responds with answers, messages, or effects in the same term-based form. Being a Prolog agent is therefore an interface and behaviour property, not a requirement about implementation language.

\medskip
\noindent\textbf{State and interaction.}
Agents are persistent and interactive, but they need not be stateful. Some are effectively stateless, while others maintain state in receive-loop parameters, control structure, or a private dynamic database. The private database is thus a capability, not a defining feature.

\medskip
\noindent\textbf{Actors and nodes.}
We treat Prolog actors and Prolog nodes uniformly as agents because they are all addressable, persistent, and interactive, and because their interactions are mediated by Prolog terms. They differ in role: actors are loci of concurrent computation; toplevels expose a conversational interface; nodes host agents and knowledge, expose web APIs, and enforce policy.

\medskip
\noindent\textbf{Private beliefs and shared knowledge.}
Each agent may combine a private workspace (including an optional private dynamic database) with shared node-provided knowledge (typically a read-only database maintained by the node owner). Private definitions can shadow shared ones, allowing local specialisation without modifying the shared base.

\medskip
\noindent\textbf{Protocols as contracts.}
Toplevel actors are special because their behaviour is governed by an explicit, observable protocol. The Prolog Toplevel Communication Protocol (PTCP) specifies how goals are submitted, how solutions are delivered incrementally, and how exceptional situations such as aborts and exits are handled. This turns REPL interaction into a first-class behavioural contract and supports later discussions of timeouts and of stateless versus session-oriented interfaces.

\medskip
\noindent The following sections refine this picture by examining actor agents and their lifecycle, then toplevel agents and PTCP, and finally nodes as agents in the Prolog Web.

%\section{Prolog agents in a nutshell}\label{sec:prolog-agents-nutshell}
%
%In this book, an \emph{agent} is understood in a deliberately modest sense: a persistent software process that participates in ongoing interaction with an external environment. This is a modelling choice rather than a claim about intelligence or autonomy.\footnote{In this spirit we follow \citet{russell2016artificial} in treating ``agent'' as an analytical lens rather than an absolute classification.}
%
%\medskip
%\noindent\textbf{Prolog agents.}
%A \emph{Prolog agent} is a process on the Prolog Web whose interactions are expressed as Prolog terms: it receives goals or messages and responds with answers, messages, or effects in the same term-based form. Being a Prolog agent is therefore an interface and behaviour property, not a requirement about implementation language.
%
%\medskip
%\noindent\textbf{State and interaction.}
%Agents are persistent and interactive, but they need not be stateful. Some are effectively stateless (an echo actor replies without remembering), while others maintain state in receive-loop parameters, control structure, or a private dynamic database. The private database is thus a capability, not a defining feature. What matters is continuity of interaction across time.
%
%\medskip
%\noindent\textbf{Actors, toplevels, and nodes as agents.}
%We treat actor processes, toplevel sessions, and nodes uniformly as agents because they are all addressable, persistent, and interactive, and because their interactions are mediated by Prolog terms. They differ in role: actors are loci of concurrent computation; toplevels expose a conversational interface; nodes host agents and knowledge, expose web APIs, and enforce policy.
%
%\medskip
%\noindent\textbf{Private beliefs and shared knowledge.}
%Each agent may combine a private workspace (including an optional private dynamic database) with shared node-provided knowledge (typically a read-only database maintained by the node owner). Private definitions can shadow shared ones, allowing local specialisation without modifying the shared base.
%
%\medskip
%\noindent\textbf{Protocols as contracts.}
%Toplevel actors are special because their behaviour is governed by an explicit, observable protocol. The Prolog Toplevel Communication Protocol (PTCP) specifies how goals are submitted, how solutions are delivered incrementally, how output and prompts may occur, and how exceptional situations such as aborts and exits are handled. This turns REPL interaction into a first-class behavioural contract and supports later discussions of timeouts and of stateless versus session-oriented interfaces.
%
%\medskip
%\noindent The following sections refine this picture by examining actor agents and their lifecycle, then toplevel agents and PTCP, and finally nodes as agents in the Prolog Web.
%
%\section{Prolog agents in a nutshell}\label{sec:prolog-agents-nutshell}
%
%The term \emph{agent} is used in many different ways across computer science and artificial intelligence, ranging from philosophical discussions of autonomy and rationality to very concrete descriptions of interactive software components. In this book, the term is used in a deliberately modest and operational sense. An agent is understood as a persistent software process that participates in ongoing interaction with an environment external to itself.
%
%This definition is intentionally broad. It does not presuppose intelligence, goal-directedness, or any particular cognitive architecture. Nor does it draw a sharp line between agents and non-agents. Rather, it serves as a modeling choice: by treating certain entities as agents, we emphasize their role as long-lived, addressable participants in structured interaction. In this sense, agenthood is a way of organizing systems rather than a claim about internal sophistication.\footnote{In other words, we follow \citet{russell2016artificial} in treating the notion of an agent as an analytical tool rather than an absolute classification.}
%
%From this perspective, many familiar software entities qualify as agents: servers that respond to requests over extended periods of time, interactive shells that maintain conversational context, background workers that receive and process messages, and controllers that coordinate access to shared resources. What distinguishes agents from short-lived computations is not intelligence, but continuity and interaction.
%
%\medskip
%\noindent\textbf{Prolog agents.}
%Within the Prolog Trinity ecosystem, a \emph{Prolog agent} is an executing software process on the Prolog Web that participates in ongoing interaction by exchanging Prolog terms with other agents. In particular, Prolog agents \emph{talk Prolog}: they receive goals, messages, or requests represented as Prolog terms, and they respond by producing answers, messages, or effects that are likewise expressed in Prolog.
%
%Importantly, this characterization does not require that a Prolog agent is written in Prolog, only that it can communicate using Prolog terms and follow the interactional conventions of the Prolog Web. Prolog agents are therefore defined by their \emph{interface and behavior}, not by their implementation language.
%
%\medskip
%\noindent\textbf{State and interaction.}
%Prolog agents are persistent and interactive, but they are not necessarily stateful. Some agents deliberately maintain no internal state beyond what is implicit in the current interaction. Echo actors are a simple example: they receive a message and immediately respond, without recording information across interactions. Such agents are stateless, yet they clearly qualify as agents by virtue of being long-lived, addressable, and continuously interactive.
%
%Other Prolog agents do maintain internal state. This state may be represented explicitly, for example by arguments threaded through a receive loop, implicitly through control flow and choice points, or explicitly in the agent's private dynamic Prolog database. The presence of such a database is therefore a \emph{capability} of Prolog agents rather than a defining requirement.
%
%What is essential is not state, but \emph{continuity of interaction}. A Prolog agent persists across multiple interactions and participates in a communication protocol that governs how messages, goals, and responses are exchanged over time. Whether or not the agent's behavior depends on accumulated internal state is an orthogonal design choice.
%
%\medskip
%\noindent\textbf{Actors, toplevels, and nodes as agents.}
%Throughout this book, several different kinds of entities are treated uniformly as Prolog agents, despite their differing responsibilities and lifetimes. Actor processes, toplevel sessions, and even entire nodes are all regarded as agents. The reason for this is not to collapse important distinctions, but to highlight a shared interaction model.
%
%Actor agents are loci of computation. They execute Prolog code concurrently, may maintain private state, and communicate by sending and receiving messages. Toplevel agents provide a conversational interface to Prolog execution, governed by an explicit protocol that makes the structure of interaction observable and programmable. Nodes, finally, act as agents at a higher level of abstraction: they host actors and databases, expose web APIs, enforce policies, and mediate access to shared resources.
%
%What these entities have in common is that they are all addressable, persistent, and interactive, and that their interactions are expressed in Prolog terms. Treating them uniformly as agents allows us to describe local and distributed interaction using a single conceptual vocabulary.
%
%\medskip
%\noindent\textbf{Private beliefs and shared knowledge.}
%Each Prolog agent has access to two distinct kinds of knowledge. The first is its private workspace, which may include a private dynamic Prolog database representing the agent's internal beliefs. This knowledge is local to the agent and inaccessible to others. The second is shared knowledge provided by the hosting node, typically in the form of a read-only shared database maintained by the node owner.
%
%This separation mirrors a common pattern in agent systems: private beliefs coexist with shared, authoritative knowledge. In the Prolog Trinity ecosystem, the distinction is enforced operationally. Private databases are process-local and cannot be modified by other agents, while shared databases are controlled by the node and exposed under clearly defined policies. Name resolution and predicate lookup are arranged so that private definitions may shadow shared ones, allowing agents to customize or extend shared knowledge locally without modifying it.
%
%\medskip
%\noindent\textbf{Agents as a unifying abstraction.}
%Referring to actors, toplevels, and nodes as agents is thus not a claim about their intelligence or autonomy, but a way of structuring the system. It provides a uniform lens through which interaction, persistence, and distribution can be described. In particular, it allows the same concepts -- conversation, addressability, lifetime, and policy -- to be applied consistently at different levels of abstraction.
%
%\medskip
%\noindent\textbf{Protocols as contracts.}
%Among the agents introduced in this chapter, toplevel actors occupy a special position. What distinguishes them is not their internal machinery, but the fact that their conversational behavior is governed by an explicit and observable protocol. The Prolog Toplevel Communication Protocol (PTCP) makes precise what clients may rely on: how goals are submitted, how solutions are delivered incrementally, how output and prompts may occur mid-computation, and how exceptional situations such as aborts and exits are handled. In this sense, PTCP turns an informal REPL convention into a first-class behavioral contract. This contract underpins many of the guarantees discussed later in the chapter, including timeout handling, API obligations, and the differences between stateless, semi-stateful, and stateful interaction.
%
%In the sections that follow, we refine this picture by examining specific kinds of Prolog agents in more detail. We begin with actor agents and their lifecycle, before turning to toplevel agents governed by an explicit communication protocol, and finally to nodes as agents that participate in the Prolog Web.


%\section{The concept of an agent}
%
%In \cite{russell2016artificial} an agent is characterized as ``anything that can be viewed as perceiving its environment through sensors and acting upon that environment through effectors.'' It does not get more abstract than that, as it covers just about everything from a simple reactive agent such as a thermostat to a complex cognitive agent such as a human.
%
%In this book, however, we shall be somewhat more concrete and care mostly about the concept of a \textit{software agent}, i.e. an agent implemented in software and running on computer hardware. The concept of software agent might be more appropriately defined as a software process capable of continuous interaction with the world external to it.  This definition of agenthood is admittedly very simple. It leaves out properties such as autonomy and goal-directedness, properties other researchers use in order to demarcate software agents from ordinary executing programs.\footnote{We agree with Russell and Norvig \citet{russell2016artificial} when they state that ``[the] notion of an agent is meant to be a tool for analyzing systems, not an absolute characterization that divides the world into agents and non-agents,'' so we do not elaborate further on this.} By our definition, any running program is an agent, as long as it is an \textit{interactive} program.\footnote{Thus, Turing machines are not agents, since they are not interactive.} 
%
%Software agents living in web-based environment form a category on its own. We shall refer to them as \textit{web agents}. A web agent, in the context of the internet and web technology, typically refers to a program or script that performs automated tasks or interacts with web services on behalf of a user or another program. %Some such agents are also known as \textit{bots} or \textit{web crawlers}. They can serve various purposes, such as web scraping, data collection, automated form submission, or even chatbots that interact with users on websites.% Web agents can be beneficial for automating repetitive tasks or gathering information from the web but can also be used for malicious purposes if not properly controlled or regulated.
%
%In the context of computer science and artificial intelligence, ``spawning an agent'' typically refers to the creation or instantiation of a new software agent from an existing one, resembling the parent-child relationship seen in biological reproduction. This process is akin to the notion of giving birth, where a ``parent agent'' generates a ``child agent.''
%
%The parent agent is responsible for initiating and managing the child agent's execution. This can involve allocating resources, defining the child agent's initial state, and establishing communication channels between the parent and child agents. The child agent inherits certain characteristics or behaviors from its parent but may also possess its own unique attributes or functionalities.
%
%This concept is frequently encountered in multi-agent systems, distributed computing environments, and parallel processing scenarios, where the dynamic creation of agents allows for increased flexibility and scalability in handling tasks and solving complex problems.
%
%\section{Prolog agents in a nutshell}\label{sec:prolog-agents-nutshell}
%
%A \textit{Prolog agent} is an executing software process on the Prolog Web which is capable of continuous interaction with other Prolog agents in its environment using Web Prolog as a dedicated agent communication language (or an ACL, to use the commonly employed acronym). In this book, we deal primarily with Prolog agents that live on the Web, or in other words, that are \textit{web agents}. Note that if taken as a definition it does not require that a Prolog agent is \textit{written} in Web Prolog, only that it can \textit{talk} Prolog. %(Further ahead in this chapter we will briefly describe a type of Prolog agent which is not implemented in Web Prolog.)%This is something that actors as well as nodes can do. 
%
%As suggested by the taxonomy in Figure~\ref{fig:prolog-agent-taxonomy}, %the Prolog Trinity ecosystem is the home to different types of Prolog agents.
%Prolog agents come in different shapes and sizes, and as suggested by the diagram, they can be arranged into a simple taxonomy. 
%
%\begin{figure}[h]
%	\centering
%	\includegraphics[width=6.5cm]{prolog-agent-taxonomy-shells}
%    \caption{A taxonomy of Prolog agents.}
%    \label{fig:prolog-agent-taxonomy}
%\end{figure}
%
%\todo{Expand with servers, supervisors expand nodes into ACTOR, ISOTOPE, etc. - perhaps also distinguish generic behaviors from specialized?}
%
%\noindent \textit{Prolog nodes} (or just nodes) serve as the runtime systems of Web Prolog, and in the context of this book they will be regarded as agents.
%\textit{Prolog actors} -- very similar to what the Erlang community refer to as processes -- are the loci of computation in the Prolog Trinity ecosystem. Prolog actors satisfy our definition of agenthood since they are processes capable of talking to other processes in their environment. As actors, they are all equipped with a mailbox from which messages received can be selected, and they can send messages to other agents. An actor can also spawn a child actor configured in all sorts of ways by means of options.
%
%There are different \textit{kinds} of actors such as \textit{toplevels} and \textit{statechart actors}, and there is a potential for other kinds of actors with other behaviors, some very specific, others generic. The ... in the diagram is a place holder for future such agents, and also covers actors such as echo actors and count actors. Chapter~\ref{sec:web-prolog} looked at many examples of such actors in action. 
%
%An explanation of what we mean by a \textit{statechart actor} will have to wait until Chapter~\ref{sec:web-prolog-statechart-behavior}. Suffice it to say that it is a way to program an actor using a (partly) visual programming language called \textit{statecharts}, invented by David Harel (who is featured among the other inventors in Chapter~\ref{sec:intro}). The relevancy to the Prolog Trinity ecosystem is that such actors can use Web Prolog as a data modelling and scripting language.
%
%If toplevels and other actors are considered agents, it follows that they represent a rudimentary form of agency, characterized by minimal complexity in decision-making and autonomy. While these agents exhibit basic properties such as autonomy, interaction, and task-specific behavior, these characteristics alone may not suffice for what might be considered a fully autonomous agent. A more robust agent, such as a ``soft robot'' interacting with and adapting to its environment, would require additional cognitive capabilities, such as the ability to form desires and intentions, as outlined in the belief-desire-intention (BDI) framework.
%
%Despite this, we remain satisfied with our simple agents, whether they function as nodes, toplevel actors, statechart actors, or other specialized entities. These agents have proven to be incredibly useful in their current form, particularly in web environments where their simplicity allows for ease of deployment and scalability. The utility of these agents is further enhanced by their ability to be composed into more complex systems, enabling the emergence of more sophisticated behaviors.
%
%Moreover, ``real'' agents can be constructed from networks or societies of these simple agents. This modular approach enables the development of more advanced agents, such as those adhering to the BDI model, by building on the cooperative interactions of basic components.

\section{More about Prolog actor agents}\label{sec:actor-agents}

Actor agents are the workhorses of the concurrent strand of the Trinity: they execute Prolog code concurrently, interact by exchanging Prolog terms, and may maintain actor-local state across repeated interactions. This section makes that operational picture more concrete. We first introduce the actor's private Prolog database and the \texttt{load\_*} options used to initialise it. We then sketch the life-cycle of a typical actor and explain how receive loops realise explicit communication protocols. Finally, we discuss the dynamic (still private) database as a non-logical but sometimes convenient mechanism for actor-local updates.


\subsection{An actor agent is equipped with a private Prolog database}\label{sec:private-db}

When running the examples in Chapter~\ref{sec:web-prolog} we were somewhat vague about the origin of the code being executed. The only clue provided was that it might reside in the node's shared database, which defines read-only data and programs accessible to any actor living on that node. While we will delve deeper into the shared database later in this chapter, we first turn our attention to the other storage area -- the one located \textit{within} the actor itself and where programs and data exclusive to that actor are stored. While all actors have a private database, it starts out empty by default. 

During the spawn operation, the (soon-to-be) parent process can initialize the private database of the (soon-to-be) child actor with clauses defining one or more predicates. Any predicate defined in the private database will take precedence over and shadow corresponding predicates in the shared database, ensuring the child actor can operate with its own customized logic and data.

Consider the following example where the \texttt{load\_text} option is passed to \texttt{spawn/3}, specifying that the source code for our count server should be loaded into the actor's private database before calling the goal in the first argument:

\begin{verbatim}
?- spawn(count_actor(0), Pid, [
       load_text("
           count_actor(Count0) :-                            
               receive({                         
                   count(From) ->
                       Count is Count0 + 1,              
                       From ! count(Count),              
                       count_actor(Count) ;
                   stop ->
                       true 
               }). 
       "),
       monitor(true)
   ]).
Pid = 45092311.
?-
\end{verbatim}

\noindent %The actor \texttt{45092311} is one hundred percent equivalent to the actor \texttt{60367387} in Section~\ref{sec:servers}, and the interaction we had with that, we can have with this.
In addition to the \texttt{load\_text} option, three other options can be used to populate the private database of an actor. The \texttt{load\_list} option loads a list of clauses or directives, the \texttt{load\_uri} option loads the content specified by a URI, and \texttt{load\_predicates} takes a list of predicate indicators and loads the clauses for the indicated predicates that are accessible by the caller. See Section~\ref{sec:actor-API} for more information. 

The clauses in an actor's private database represent its unique beliefs and skills, distinguishing them from the shared beliefs and skills accessible to other actor agents on the same node. In this sense, the private database can be seen as the core of the actor's individuality, encapsulating the knowledge and capabilities that are exclusive to it.

The \texttt{load\_*} options make a kind of sharing of beliefs and skills between a parent process and a child process possible. The child is ``born'' with these beliefs and skills, they are ``innate.'' It is up to the parent actor to determine exactly what should be passed on to the child actor. %We can think of this as ``learning by injection''.

At this point, and practically speaking, the actor's private database and the \texttt{load\_*} options may not seem all that useful. After all, it seems we could just as well have stored the implementation of \texttt{count\_actor/1} in the node's shared database. As we shall see, they are more useful in the distributed case, for shipping predicates from a client to a node, or from a node to a client, and execute them there.

%\subsubsection*{Shared data and private beliefs}\label{sec:shared-vs-private}
%
%The node-resident programs and data can be seen as ``shared knowledge'' serving as a background to the ``private beliefs'' of any Prolog actor that currently happens to live on the node in question. However, note that the shared knowledge in reality can be quite extensive, as the resident program on this node may be linked -- and linked rather seamlessly -- to resident programs and data tied to other nodes. 


\subsection{The life-cycle of a typical Prolog actor agent}\label{sec:actor-lifecycle}

\begin{quotation}
\noindent We have the black boxes that do things, and we have the messages in between. What goes on inside the black boxes does not really matter; what takes place there is abstracted over.
\vspace{-0.1cm}
\flushright  \textit{Joe Armstrong}
\vspace{0.1cm}
\end{quotation}

\noindent The diagram in Figure~\ref{fig:life-cycle} illustrates the life-cycle of a typical Prolog actor. 

\begin{figure}[h]
\centering
	\includegraphics[width=7.9cm]{life-cycle-kombi}
    \caption{Lifecycle of a Prolog actor: spawn, interact, terminate, and clean up}
    \label{fig:life-cycle}
\end{figure}

When an actor is spawned, it first passes through an initialization phase where it is configured in a way that will suit its intended purpose. This includes specifying values for options such as \texttt{monitor}, \texttt{link} as well as the \texttt{load\_*} options.

For example, when the count actor in the previous sections was spawned, the source code specified by the \texttt{load\_text} option was loaded into the actor's private Prolog database and the \texttt{monitor} options instructed the actor to send a \texttt{down} message to the shell just before terminating. The call to \texttt{spawn/3} succeeded, binding the variable in the second argument to a pid. By this, the initialization phase is completed, the actor is now in its \textit{initial state}, and the loop that implements the actor's \textit{communication protocol} is entered. Typically, processing stops at a call to \texttt{receive/1-2} where it blocks, waiting for a \texttt{count} message or a \texttt{stop} message to appear in its mailbox.

%At this point, a client that knows the pid of an actor process (or its name if it has one) and knows its communication protocol (or has access to predicates that hides its details), is invited to send it messages. In this way, the behavior of the actor can be controlled and its functionality harnessed.

The notion of a communication protocol is central to Erlang-style concurrent programming. The loop implements the communication protocol that clients that want to interact with the actor must follow. The interaction is driven by coordinating send and receive operations, with the protocol defining the rules. The sending process transmits messages in the specified format, and the receiving process captures and interprets them. They must adhere to the syntax, semantics and timing constraints of the communication protocol:

\begin{itemize}

\item \textbf{Syntax}: This refers to the structure and format of the data exchanged between processes. It dictates how data is packaged, including the organization of various data elements.

\item \textbf{Semantics}: This involves the meaning and interpretation of the exchanged data. It defines the rules for understanding the content of the message, what actions should be triggered upon receiving certain data, and how to handle errors or unexpected data.

\item \textbf{Timing}: This concerns the timing aspects of the communication, such as when data is sent, the pace of data exchange, and managing the order and sequence of messages so that both processes can coordinate despite delay, reordering, and failure.

\end{itemize}

\noindent The Web Prolog receive operator, while blocking the process at the point of a \texttt{receive/1-2} call, allows for the implementation of very flexible communication protocols through mechanisms like selective pattern matching and timeouts. The selective pattern matching allows the process to wait for a particular message but can also defer out-of-order messages for later processing. By specifying a timeout to avoid the receive call being blocked indefinitely, alternative actions can be taken if an expected message does not arrive in time. This helps reduce the chance of breakdowns and allows for more adaptive communication.

Whenever the actor runs a \texttt{receive/1-2} call, a message is selected from the mailbox and processed sequentially. The result often depends on the state of the process, and as part of processing, a new state for the actor may be computed. This will be the current state for processing the next message selected.

Some actors are stateless (e.g.\ the echo actor), and those that are stateful represent their state and perform state transitions in a variety of ways: as values threaded through the arguments of a receive loop (for example an integer updated by calls to \texttt{is/2} in the count actor, or a list of atoms in the fridge simulation), as the explicit states and transitions of a state machine, or implicitly through the choice points created by a failure-driven loop. These techniques all encode state in the control flow or data flow of the program itself. In addition to these, Web Prolog offers a more direct -- though non-logical -- mechanism for maintaining actor-local state: a private dynamic Prolog database.

\subsection{The dynamic (and still private) Prolog database}\label{sec:dynamic-private-database}

\noindent As long as the actor process is alive, it is allowed to update its private database using predicates such as \texttt{assert/1}, \texttt{retract/1} and \texttt{retractall/1}. Following the ISO standard, predicates that are modified this way need to be declared using the \texttt{dynamic/1} directive. Updates only affect that actor's private database: they are not visible to other actors running on the same node.

Because the dynamic database is strictly actor-local, it neither introduces shared-state races nor provides a mechanism for shared-state inter-process communication.

For example, although the following goal is permitted, it is of limited use: the clause \texttt{foo(a)} exists only inside the spawned actor and disappears when that actor terminates.

\begin{verbatim}
?- spawn(assert(foo(a)), Pid).
Pid = 32861299.
?-
\end{verbatim}

\noindent The dynamic database provides another way to transfer from one state to the next. This is often a bad idea.

\medskip
\noindent
Throughout the Prolog Trinity ecosystem, private state is treated as volatile and local to an agent, whereas shared data is governed by explicit node-level policies and exposed under controlled conditions.


%\subsection{The dynamic (and still private) Prolog database}\label{sec:dynamic-private-database}
%
%\noindent As long as the actor process is alive, it is allowed to update its private database using predicates such as \texttt{assert/1}, \texttt{retract/1} and \texttt{retractall/1}. Following the ISO standard, predicates that are modified this way need to be declared using the \texttt{dynamic/1} directive. Updates will only affect that actor's private database, not actors running elsewhere, and not even actors running on the same node. 
%
%Because the dynamic database is strictly actor-local, it neither introduces shared-state races nor provides a mechanism for inter-process communication.
%
%%This means that problems caused by two or more processes trying to update the same database simultaneously cannot arise. It also means that since database updates performed by one process is not seen by other processes, assert and retract cannot be used for inter-process communication. Web Prolog adheres to the idea that processes should ``communicate to share memory, rather than share memory to communicate.''\footnote{A mantra attributed to CSP.}
%
%Therefore, although the following goal is permitted, it is quite meaningless since the clause \texttt{foo(a)} disappears as the goal has run and the spawned process has terminated.
%
%\begin{verbatim}
%?- spawn(assert(foo(a)), Pid).
%Pid = 32861299.
%?-
%\end{verbatim}
%
%\noindent The dynamic database provides another way to transfer from one state to the next. This is often a bad idea. 
%
%\medskip
%\noindent
%Throughout the Prolog Trinity ecosystem, private state is treated as volatile and local to an agent, whereas shared data is governed by explicit node-level policies and exposed under controlled conditions.

\subsubsection{More about the dynamic database}\label{sec:dynamic-database}

The dynamic database is a non-logical extension to Prolog, and its use for managing state is often viewed with suspicion by Prolog programmers. The reason is well known: changes to the database are not reverted on backtracking, which makes dynamic predicates behave like mutable global variables and exposes programs to the usual problems associated with such entities.

At first glance, this raises the question of whether the dynamic database could lead to problems in Web Prolog, especially in a concurrent setting. The short answer is no. Although Erlang does not offer a Prolog-style dynamic database, it provides a closely related mechanism in the form of the \textit{process dictionary}. Each Erlang process maintains a private key--value store that cannot be accessed or modified by other processes. As with the dynamic database, its use is generally discouraged, since it undermines referential transparency and makes programs harder to reason about and debug.

The similarity is hard to miss. This suggests that the dynamic database in a Web Prolog actor is not inherently more dangerous than the process dictionary in Erlang. At the same time, the dynamic database is arguably more expressive and convenient, which also means that programmers may be tempted to rely on it more often than is advisable.

While both the dynamic database in Web Prolog actors and the Erlang process dictionary are strictly private, Erlang also provides mechanisms for mutable \emph{shared} state. These include Erlang Term Storage (ETS), which offers in-memory tables with configurable access rights, and Mnesia, a distributed database supporting transactions, persistence, and replication. More recently, Erlang has also introduced \texttt{persistent\_term}, which provides efficient read-only access to immutable global data, at the cost of expensive updates.

This raises a broader design question for Web Prolog: should it offer forms of dynamic storage beyond what can be represented by predicate arguments or asserted in an actor's private database? Given Web Prolog's focus on the Web, a transactional RDF store would be a natural candidate, and such a solution is readily available in systems like SWI-Prolog. A Linda-style blackboard is another possible direction. For the time being, however, Web Prolog deliberately restricts dynamic state to the private scope of individual agents, leaving shared mutable state to be addressed by explicit, higher-level mechanisms.


%\subsubsection{More about the dynamic database - OLD}\label{sec:dynamic-database}
%
%\noindent The dynamic database is a non-logical extension to Prolog and the use of it for managing state is often frowned upon by Prolog programmers. The reason is that changes to the database are not reverted on backtracking and thus clauses of dynamic predicates are like mutable global variables and therefore subject to the usual problems with such entities. 
%
%A Prolog-style dynamic database is not something that Erlang offers, and we need to determine if it could lead to problems in Web Prolog. The answer is ``no" since Erlang does in fact have something which plays the same role as the dynamic database. In Erlang, each process has a local store called the \textit{process dictionary} and a handful of functions for manipulating it are built into the language. Process dictionaries are private and cannot be modified by any other process. Erlang programmers are advised against using them since they destroy referential transparency and make debugging difficult.
%
%One cannot help noticing a similarity here. This leads us to believe that the dynamic database in Web Prolog is not more dangerous than the process dictionary in Erlang. However, it is arguably also more powerful and useful than the process dictionary, which probably means that programmers are tempted to use it more often than is recommended.
%
%Whereas the dynamic database in a Web Prolog actor process and the process dictionary of an Erlang process are private, Erlang also supports a built-in term storage known as \textit{ETS} and a distributed database named \textit{Mnesia}. They are not only mutable but also \textit{shared} between processes. 
%
%The question whether or not Web Prolog should allow for some kind of dynamic memory over and above what can be passed around in the arguments of predicates or asserted in the private dynamic database is interesting. In line with the focus of Web Prolog on the web domain, a solution based on a transactional RDF database might be appropriate and is easily implemented in (at least) SWI-Prolog, which has a library for it.
%
%A Linda blackboard might be another alternative.
%
%---
%
%In Erlang, the process dictionary provides a means of storing key-value pairs within an individual process. Each Erlang process maintains its own isolated dictionary, which is inaccessible to other processes, reinforcing the language's fundamental principle of process isolation. This design ensures that shared mutable state does not exist by default, thereby facilitating Erlang's lightweight concurrency model and fault tolerance. Although the process dictionary can be convenient for storing process-specific data without explicit message passing, its use is generally discouraged outside of specific performance-critical scenarios due to its implicit nature, which can lead to code that is harder to reason about and debug.
%
%For cases where global state or shared access to data is required, Erlang provides alternative mechanisms that align with its concurrency model. One approach is to use a registered process, which acts as a global entity that maintains state and responds to requests from other processes via message passing. This approach ensures controlled access to shared data, preserving Erlang's principles of explicit communication and process encapsulation, though it may introduce a bottleneck depending on usage patterns. A more efficient option for concurrent access is Erlang Term Storage (ETS), which provides an in-memory key-value store that multiple processes can read from and write to while still maintaining process isolation to a degree. ETS tables can be configured with different access levels, allowing processes to share data without directly violating the shared-nothing principle.
%
%In scenarios where durability and transactional consistency are required, Mnesia, Erlang's distributed database, provides a more advanced solution by enabling persistent storage and replication across nodes. Mnesia supports transactions and complex queries, making it well-suited for applications requiring strong consistency guarantees. Another modern approach is the \texttt{persistent\_term} module, which allows for storing immutable global data with highly efficient read access. Although \texttt{persistent\_term} is optimized for read-heavy workloads, updates are costly and should be used sparingly. Each of these mechanisms offers a structured way to handle shared state without compromising Erlang's core design principles, in contrast to the process dictionary, which remains an inherently local storage facility confined to the process that owns it.

% \section{Web Prolog as a scripting language for the Web}\label{sec:web-prolog-inspired-by-javascript}

% \begin{quotation}
% \noindent Scripting languages are a lot like obscenity. I can't define it, but I'll know it when I see it.
% \flushright\textit{Larry Wall}
% \vspace{0.3cm}
% \end{quotation}

% \noindent Over the years, the Web has become a key delivery platform for a variety of sophisticated interactive applications in just about any conceivable domain. As a consequence, JavaScript, more or less the only game in town for programming the front-end of a web application, has become among the most commonly used programming languages on Earth. It appears that even back-end developers are more likely to use it than any other language.\footnote{\url{http://stackoverflow.com/research/developer-survey-2016\#most-popular-technologies-per-occupation}} Indeed, JavaScript must be regarded as \textit{the} web programming language of our times. 

% JavaScript started out as a very small, highly domain-specific scripting language that was limited to running within a web browser to dynamically modify the web page being shown, but has over time evolved into a widely portable general-purpose programming language. Modern JavaScript is a high-level, dynamically typed, interpreted multi-paradigm programming language with several essential features that make it a versatile tool for web development. JavaScript\textquotesingle s syntax supports various programming paradigms, including imperative, functional, and object-oriented programming styles. With features such as callbacks, promises, and async/await, JavaScript supports asynchronous programming, enabling non-blocking operations, especially useful in web applications. JavaScript\textquotesingle s native format for data interchange, JSON, is lightweight and widely used for data transmission. JavaScript runs on almost all modern web browsers and platforms, making it universally applicable for web development. It easily interfaces with numerous web APIs (Application Programming Interfaces) for tasks like accessing the Document Object Model (DOM), making HTTP requests and handling events. This makes it an excellent tool for connecting disparate systems in web applications.

% With the advent of Node.js,\footnote{\url{https://en.wikipedia.org/wiki/Node.js}} JavaScript extended its reach to server-side development. This allowed for a unified language experience across both client and server sides, simplifying the development process by allowing the same language to be used throughout the entire stack. Not having to learn more than one language in order to become a so called ``full-stack developer,'' and not having to switch languages when changing from working on the server-side to working on the client-side are important advantages. Also, Node.js is frequently used to build and connect microservices. Its non-blocking I/O model and event-driven architecture make it suitable for building scalable and efficient network applications.

% It is clear that as a relational logic programming language with the essential and important features listed above, Prolog is very different from JavaScript in terms of paradigms supported. The addition of Erlang-style concurrency and asynchronous communication does not really narrow the gap, as those are features that work differently in Erlang. What matters here is the \textit{role} that JavaScript plays on the Web. We want Web Prolog to play a similar role in the Prolog Trinity ecosystem as JavaScript does in the JavaScript ecosystem, just as clear-cut.

% ---

% Paradigms alone hardly capture all there is to a practical programming language. Languages may also be related by having a special purpose in common, or by constraints imposed upon them by a particular area of application. They are related by family resemblances rather than anything else, and connected by a series of overlapping similarities, where no one feature is common to all languages. 

\section{Prolog shells and other toplevel actors}\label{sec:toplevel}

As we saw already in Chapter~\ref{sec:web-prolog}, calling \texttt{self/1} in a terminal binds its argument to the pid of the currently running shell process:

\begin{verbatim}
?- self(Pid).
Pid = 72097632.
?-
\end{verbatim}

\noindent Here, the actor with the pid \texttt{72097632} is a shell. Shells are toplevels, and therefore also agents. Shells handle all the things that toplevels handle, but also add a number of useful things for when we are talking to Prolog over a terminal. This includes I/O (read and write) and utilities such as \texttt{flush/0} and the dollar notation. Such features are supplied by the \textit{node controller} rather than a toplevel alone. (We will say more about the node controller further ahead in this chapter.)

Here is how we instruct the shell to spawn a new toplevel:

\begin{verbatim}
?- toplevel_spawn(Pid).
Pid = 16226587.
?-
\end{verbatim}

\noindent Note that \texttt{16226587} is a toplevel actor but not a shell. As with every other Prolog actor, it comes with a mailbox, a private Prolog database, the ability to send messages to other actors, as well as the ability to create other actors.

The feature that distinguishes a toplevel from other actors is that it comes with a standardized built-in special-purpose communication protocol. 
%, the Prolog Toplevel Communication Protocol.

\begin{figure}[h]

\centering
	\includegraphics[width=6cm]{a-toplevel-actor}
    \caption{The anatomy of a Prolog toplevel actor.}
    \label{fig:a-toplevel-actor}
\end{figure}

\subsection{A Prolog toplevel is an actor with a built-in protocol}

A programmer firing up a traditional Prolog system is likely met with a query prompt. In the literature, this is usually referred to as the \textit{toplevel}. The reason we refer to it a \textit{shell} is because we want to use the term \textit{toplevel} to refer to toplevels that are not shells. 

In traditional Prolog, a program cannot \emph{internally} create a toplevel, pose queries and request solutions on demand, but this is something that Web Prolog allows. In traditional Prolog the shell is \textit{lazy} in the sense that new solutions to a query are only computed on demand. As we shall see, this is how the Web Prolog toplevel actor works too.

In Web Prolog, a toplevel actor is a programming abstraction modelled on the interactive toplevel of Prolog. A toplevel actor is like a first-class interactive Prolog toplevel, accessible from Web Prolog as well as from other programming languages such as JavaScript. We can also think of it as an \textit{encapsulated Prolog session}, an abstraction aiming at making Prolog programmers feel right at home.

A toplevel is a kind of actor, and what distinguishes it from other kinds of actors is the \textit{protocol} it follows when it communicates, i.e. the kind of messages it listens for, the kind of messages it sends and in what order, and the behavior this gives rise to. 

The protocol must not only allow a client to submit queries and a toplevel to respond with answers, it must also allow the toplevel to  prompt for input or produce output at any time, in an order and with a content as dictated by the program that it runs. All toplevels follow this protocol. The terminal adheres to it as well, and even a human user of a terminal talking to a shell must adapt to it in order to have a successful interaction.

The design of the toplevel actor in Web Prolog is in fact very much inspired by the informal communication protocol that we as programmers adhere to when we invoke a Prolog shell from our OS prompt, load a program, submit a query, are presented with a solution (or a failure or an error), type a semicolon in order to ask for more solutions, or hit return to stop. These are ``conversational moves'' that Prolog understands. There are even more such moves, because after having run one query to completion, the programmer can choose to submit another one, and so on. The session does not end until the programmer decides to terminate it. There are only a few moves a client can successfully make when the protocol is in a particular state, and the possibilities can easily be described, by a state machine for example, as will be shown in the next section.

\subsection{The Prolog Toplevel Communication Protocol}\label{sec:pcp-protocol}

\noindent The diagram in Figure~\ref{fig:statechart-2} is a statechart specifying the Prolog Toplevel Communication Protocol (PTCP) -- a protocol for the communication between a client and a toplevel actor. The client can be any process (including another actor or (say) a JavaScript process) capable of sending the messages and signals in bold to the toplevel. The toplevel actor is responsible for returning the messages with a leading / back to the client. The use of a statechart allows us to show that no matter the current state of the protocol, \textbf{abort} will always take it to the state from which a new goal can be called and \textbf{exit} will always terminate the toplevel process.  %\footnote{The use of a statechart allows us to show that no matter the current state of the protocol, \textbf{abort} will always take it to the state from which a new query can be asked and \textbf{exit} will always terminate the toplevel process.}

\begin{figure}[h]
    \centering
	\includegraphics[width=9cm]{PLAP-statechart-3-states.pdf}
 %   \caption{Statechart specifying the PTCP for a successful conversation with a toplevel. The transitions are labeled with \textit{message types}. Types in bold are sent from the client to the toplevel, whereas message types with a leading / goes in the opposite direction, from the toplevel to the client.}
    \caption{Statechart for the Prolog Toplevel Communication Protocol (PTCP): client moves and toplevel responses.}
    \label{fig:statechart-2}
\end{figure}

A process is able to \textit{spawn} a toplevel. After having done so, the process becomes the parent of the toplevel and can start communicating with it according to the protocol. Initially, the protocol is in state \textbf{s1}, where the toplevel rests idle, waiting for a \textbf{call} message containing a goal. When such a message arrives, the protocol transitions to state \textbf{s2}, where the toplevel is doing actual work. The protocol will remain in state \textbf{s2} until some work is done and the toplevel sends a message indicating either \texttt{/success}, \texttt{/failure}, \texttt{/error}, or (if the process is monitored) \texttt{/down} to the client. On the event of the toplevel sending \texttt{/failure} or \texttt{/error}, the protocol will transition back to state \textbf{s1}. This will be the case also for a \texttt{/success} message that indicates no more solutions to the query can be found (marked with false in the chart). However, if a \texttt{/success} message indicates more solutions may exist (marked with true in the chart), the protocol transitions to state \textbf{s3}, where it will wait for a message \textbf{next}, \textbf{stop}, \textbf{abort} or \textbf{exit} to arrive from the client. If the message is \textbf{stop} or \textbf{abort} the protocol will transition back to state \textbf{s1}, if it is \textbf{next} it will transition to state \textbf{s2} and trigger the search for more solutions, and finally, if it is \textbf{exit}, it will (if the process is monitored) force the toplevel to send a message \texttt{/down} back to the client and then to terminate. Any other messages will be deferred to the code being executed by the toplevel.

Web Prolog comes with built-in predicates which allow a client to spawn a toplevel actor and interact with it during its life-cycle:

\begin{verbatim}
toplevel_spawn/1-2    toplevel_call/2-3    toplevel_next/1-2
toplevel_stop/1       toplevel_abort/1 
\end{verbatim}

\noindent As witnessed by some of the examples in Chapter~\ref{sec:web-prolog}, such predicates can be defined in terms of the more primitive \texttt{spawn/1-3}, \texttt{!/2} and \texttt{receive/1-2} predicates. The predicates proposed here come with a number of important additional options, the use of which will be demonstrated in the sections that follow. Appendix \ref{sec:manual} contains an excerpt from the (draft) manual which covers most of the details of our proposal.

\subsubsection*{Toplevel actors are stateful}\label{sec:toplevel-stateful}

Clearly, a Prolog toplevel is a kind of server (in the sense of Erlang -- see Section~\ref{sec:servers}). Also, note that a toplevel, even when not explicitly threading any state, nor using the dynamic database, must still be considered stateful. Rather than using an explicit data structure for holding the state, it is the underlying Prolog process \textit{as such} that holds it, most clearly shown in the way the toplevel ``remembers'' its history and how its behavior is influenced by this, enabling it to react appropriately when a client requests the next solution to a query.


\subsection{Signatures and options for \texttt{toplevel\_*} predicates}

\subsubsection*{\texttt{toplevel\_spawn(-Pid, +Options?)}}

% Requires: \usepackage{tabularx}

\begin{tabularx}{11cm}{|X|X|X|}
\hline
\textbf{Option} & \textbf{Values} & \textbf{Default} \\
\hline
\texttt{session} & \texttt{true}, \texttt{false} & \texttt{false} \\
\texttt{target} & Pid & caller's pid \\
\hline
\end{tabularx}

\subsubsection*{\texttt{toplevel\_call(+Pid, +Goal, +Options?)}}

\begin{tabularx}{11cm}{|X|X|X|}
\hline
\textbf{Option} & \textbf{Values} & \textbf{Default} \\
\hline
\texttt{template} & term & Goal \\
\texttt{offset} & integer $\geq$ 0 & \texttt{0} \\
\texttt{limit} & integer $>$ 0 & all \\
\texttt{once} & \texttt{true}, \texttt{false} & \texttt{false} \\
\texttt{target} & Pid & from \texttt{toplevel\_spawn} \\
\hline
\end{tabularx}

\subsubsection*{\texttt{toplevel\_next(+Pid, +Options?)}}

\begin{tabularx}{11cm}{|X|X|X|}
\hline
\textbf{Option} & \textbf{Values} & \textbf{Default} \\
\hline
\texttt{limit} & integer $>$ 0 & from \texttt{toplevel\_call} \\
\texttt{target} & Pid & from \texttt{toplevel\_call} \\
\hline
\end{tabularx}

\subsection{Shell talking to a Prolog toplevel}

\noindent Below, we show an example of how to create and interact with a toplevel process from a shell. We start by spawning a new toplevel, using \texttt{toplevel\_spawn/2} with the \texttt{monitor} option to instruct the toplevel to send us a \texttt{down} message when the process eventually terminates. The \texttt{load\_list} option is used to populate the private Prolog database belonging to the toplevel actor with two simple unit clauses \texttt{p(a)} and \texttt{p(b)}:

\begin{verbatim}
?- toplevel_spawn(Pid, [
       session(true),
       monitor(true),
       load_list([p(a),p(b)])
   ]).
Pid = 74122981.
?-
\end{verbatim}

\noindent With this, the toplevel actor is initiated, and with the PTCP now in  state \textbf{s1}, it is ready to accept messages and signals sent by the other \texttt{toplevel\_*} predicates.

The \texttt{monitor} and \texttt{load\_list} options are inherited from \texttt{spawn/1-3}. However, \texttt{toplevel\_spawn/2} offers a couple of additional options that can be used to specify the behavior of toplevel actor at creation time. In the above example the \texttt{session} option is used to configure the toplevel to run a multi-goal session rather than exit after having run one goal to completion. 


\subsubsection{Making the toplevel call a goal}\label{sec:toplevel-call-goal}

Let us see what happens if we call \texttt{toplevel\_call/2} with the default values for options:

\begin{verbatim}
?- toplevel_call($Pid, p(X)).
true.
?- flush.
Shell got success(74122981,[p(a),p(b)],false)
true.
?-
\end{verbatim}

\noindent The answer is returned to the mailbox of the calling process in the form of an \textit{answer term} with three arguments. The functor of the term represents its \textit{type}. In this case, it shows that the goal succeeded. (In other cases it might indicate a failure or an error.) The first argument of the success term is the pid, and the list in the second argument represents the two solutions that were computed. The value \texttt{false} in the third argument of the term indicates that no more solutions exist. 

Since the \texttt{session} option was set to true, the PTCP is now, after a brief visit to state \textbf{s2} for the execution of the goal, back in state \textbf{s1} and is prepared to accept another use of \texttt{toplevel\_call/2-3}.

\subsubsection{Using the template option}\label{sec:toplevel-template}

Below, \texttt{toplevel\_call/3} is called with the goal \texttt{p(X)} and with the \texttt{template} option set to the variable \texttt{X}:

\begin{verbatim}
?- toplevel_call($Pid, p(X), [
       template(X)
   ]).
true.
?- flush.
Shell got success(74122981,[a,b],false)
true.
?-
\end{verbatim}

\noindent Note how the value of the \texttt{template} option determined the form of the list of solutions in the second argument of the answer term. The relation to the Prolog built-in standard predicate \texttt{findall/3} is evident. %, with its template, goal and list of solutions, is evident -- the difference is that w

\subsubsection{Using the \texttt{offset} and \texttt{limit}  options}\label{sec:toplevel-paging}

When querying a relational database using SQL or an RDF dataset using SPARQL, the use of the \texttt{OFFSET} and \texttt{LIMIT} keywords serves to control the subset of results that are returned by a query. In Web Prolog they can be used as options for the configuration of calls to \texttt{toplevel\_call/3}. The \texttt{limit} option specifies the maximum number of solutions to return, while the \texttt{offset} option specifies the number of solutions to skip before starting to return solutions. Here is what happens if \texttt{limit} is set to \texttt{1}:

\begin{verbatim}
?- toplevel_call($Pid, p(X), [
       template(X),
       limit(1)
   ]).
true.
?- flush.
Shell got success(74122981,[a],true)
true.
?-
\end{verbatim}

\noindent The value \texttt{true} in the \texttt{success} term means that the state machine is now in state \textbf{s3}, waiting for a message \textbf{next} or \textbf{stop} to appear. Calling \texttt{toplevel\_next/1} produces the next solution:

\begin{verbatim}
?- toplevel_next($Pid).
true.
?- flush.
Shell got success(74122981,[b],false)
true.
?-
\end{verbatim}

\noindent \texttt{toplevel\_call/3} supports the \texttt{offset} option. Together with the \texttt{limit} option it allows the call to pick out an arbitrary slice of solutions to a goal, even when the number of solutions is infinite:\footnote{It may seem odd and not all that useful for \texttt{toplevel\_call/2-3} to support the \texttt{offset} option, but it turns out to be important when implementing the stateless HTTP protocol.} 

\begin{verbatim}
?- toplevel_call($Pid, between(1,infinite,I), [
       offset(100),
       template(I),
       limit(3)
   ]).
true.
?- flush.
Shell got success(74122981,[101,102,103],true)
true.
?-
\end{verbatim}

\noindent Calling \texttt{toplevel\_next/1} produces three more solutions:

\begin{verbatim}
?- toplevel_next($Pid).
true.
?- flush.
Shell got success(74122981,[104,105,106],true)
true.
?-
\end{verbatim}

\noindent Note that the value of the \texttt{limit} option set in the call to \texttt{toplevel\_call/3} was still in effect. However, \texttt{toplevel\_next/2} accepts the \texttt{limit} option too, which allows us to increase or decrease the limit on the number of solutions returned:

\begin{verbatim}
?- toplevel_next($Pid, [ 
       limit(5)
   ]).
true.
?- flush.
Shell got success(74122981,[107,108,109,110,111],true)
true.
?-
\end{verbatim}

\noindent Calling \texttt{toplevel\_stop/1} will stop the execution of the goal:

\begin{verbatim}
?- toplevel_stop($Pid).
true.
?-
\end{verbatim}

\medskip
\noindent These options are particularly useful in several scenarios: pagination of results in web applications, where it is impractical to retrieve and display all records at once; implementing infinite scrolling interfaces that load data dynamically as the user scrolls; handling goals with very many or infinitely many solutions without overwhelming the client; and reducing the number of network roundtrips on the Prolog Web by fetching solutions in batches. A special case is when \texttt{limit} is set to \texttt{1}, which makes it straightforward to implement a Prolog-style REPL in JavaScript, where each press of \texttt{;} triggers a new request for the next solution.

\subsubsection{Using the \texttt{once} option}\label{sec:toplevel-once}

The \texttt{once} option only makes sense in combination with the \texttt{limit} option. Default is \texttt{false} but when set to \texttt{true}, the list of solutions in a \texttt{success} answer term are the only solutions that the call will produce (as indicated by \texttt{false} in the third argument of the answer term):

\begin{verbatim}
?- toplevel_call($Pid, between(1,infinite,I), [
       offset(100),
       template(I),
       limit(3),
       once(true)
   ]).
true.
?- flush.
Shell got success(23091243,[101,102,103],false)
true.
?-
\end{verbatim}


\subsubsection{A toplevel can do a kind of I/O}\label{sec:toplevel-io}

\noindent Whenever a toplevel is in state \textbf{s2} and doing some real work, it is able to send messages to its parent using the ordinary send operator. However, a better idea can often be to use \texttt{output/1}:

\begin{verbatim}
?- toplevel_call($Pid, output(hello)).
true.
?- flush.
Shell got output(74122981,hello)
Shell got success(74122981,[output(hello)],false)
true.
?-
\end{verbatim}

\noindent Input can be collected by calling \texttt{input/2}, which sends a \texttt{prompt} message to the client, which in turn can respond by calling \texttt{respond/2}:

\begin{verbatim}
?- toplevel_call($Pid, input('Input', X)),
   receive({Answer -> true}).
Answer = prompt(74122981,'Input').
?- respond($Pid, hello), 
   receive({Answer -> true}).
Answer = success(74122981,[input('Input',hello)],false).
?- 
\end{verbatim}

\subsubsection{Aborting a nonterminating goal}\label{sec:toplevel-abort}

\noindent Our toplevel actor is still not dead so let us see what happens when we ask the toplevel to first update its private Prolog database with a silly recursive clause for a predicate \texttt{p/0} and then call it:

\begin{verbatim}
?- toplevel_call($Pid, assert((p :- p))),
   receive({Answer -> true}).
Answer = success(74122981,[assert((p:-p))],false)
?- toplevel_call($Pid, p).
true.
?-
\end{verbatim}

\noindent Although nothing is shown in the terminal, it is clear that the toplevel is now just wasting CPU cycles to no avail. Fortunately, a nonterminating goal can be aborted by calling \texttt{toplevel\_abort/1}:

\begin{verbatim}
?- toplevel_abort($Pid).
true.
?-
\end{verbatim}

\noindent With this, the PTCP is back in state \textbf{s1}.

\subsubsection{Exiting the toplevel}\label{sec:toplevel-exit}

\noindent When we are done talking to the toplevel we can kill it. As the \texttt{monitor} option was set to \texttt{true}, we should expect to receive a \texttt{down} message:

\begin{verbatim}
?- exit($Pid, goodbye),
   receive({Answer -> true}).
Answer = down(74122981,74122981,goodbye).
?-
\end{verbatim}

\subsubsection{Toplevels and the message deferring mechanism}\label{sec:deferring}

In the following example, a toplevel is spawned, and then \texttt{toplevel\_next/1} is called. The protocol in state \textbf{s1}, which is obviously not in a state where it can react on the \textbf{next} message (see Figure~\ref{fig:statechart-2}). The message is therefore deferred and it is not until \texttt{toplevel\_call/3} is called and the protocol changes states that it has an effect.

\begin{verbatim}
?- toplevel_spawn(Pid).
Pid = 78340943.
?- toplevel_next($Pid).
true.
?- flush.
true.
?- toplevel_call($Pid, member(X,[a,b,c]), [
       limit(1),
       template(X)
   ]).
true.
?- flush.
Shell got success(78340943, [a], true)
Shell got success(78340943, [b], true)
true.
?- 
\end{verbatim}

\noindent Note that messages sent to a toplevel will often be handled in the right order even if they arrive in the ``wrong'' order (e.g. \textbf{next} before \textbf{call}). This is due to the selective receive which defers their handling until the PTCP permits it. The messages \textbf{abort} and \textbf{exit}, however, will never be deferred. This is guaranteed by the fact that \textbf{abort} and \textbf{exit} are valid transitions from any state in the protocol (see Figure~\ref{fig:statechart-2}). (They are actually signals rather than messages).

One way to think of this is in terms of the so called \textit{robustness principle}: ``Be conservative in what you send, be liberal in what you accept.''\footnote{\url{https://en.wikipedia.org/wiki/Robustness\_principle}} Due to the deferring behavior a toplevel is liberal in this way, but, as implied by the principle, clients are advised not to rely on this behavior.

\subsubsection{Redirecting answers}\label{sec:toplevel-redirect}

The \texttt{target} option can be used to redirect the answer terms produced by the toplevel to any actor of choice. By default, it is set to the pid of the parent but allows us to instruct a toplevel actor to send answer terms to a different destination. To demonstrate how this works, we first create a simple actor that can serve as a target:

\begin{verbatim}
?- spawn(( repeat,
           receive({
               Msg -> 
                   format("Received ~p~n", [Msg]), 
                   fail
           })
         ), Pid0).
Pid0 = 98380209.
?-
\end{verbatim}

\noindent The pid of that actor can now be used as the value of the target option:

\begin{verbatim}
?- toplevel_spawn(Pid, [
       target($Pid0)
   ]).
Pid = 97919106.
?-
\end{verbatim}

\noindent When \texttt{toplevel\_call/3} is called, the \texttt{success} answer term is printed by the designated target actor:

\begin{verbatim}
?- toplevel_call($Pid, between(1,1000,N), [
       template(N),
       limit(5)
   ]).
true.
Received success(97919106,[1,2,3,4,5],true)
?-
\end{verbatim}

\noindent So, \texttt{98380209} rather than the shell received the message from \texttt{97919106}.

Calling \texttt{flush/0} shows that the mailbox belonging to the shell is empty:

\begin{verbatim}
?- flush.
true.
?-
\end{verbatim}

\noindent The \texttt{target} option is supported by \texttt{toplevel\_spawn/2}, \texttt{toplevel\_call/3} and \texttt{toplevel\_next/2}. It is not supported by the other \texttt{toplevel\_*} predicates as they do not produce answer terms. It can be used in calls to \texttt{output/2}. See Appendix~\ref{sec:manual} for more details.

\subsection{Programming with toplevel actors}\label{sec:programming-with-toplevels}

The toplevel actor is an important kind of agent in the Prolog Trinity ecosystem. However, the reasons for its importance will have to be deferred to later chapters, and this section only offers a few exercises to build familiarity with it.

\subsubsection{A synchronous predicate API to toplevels}\label{sec:toplevel-sync-api}

The communication between a client and a toplevel illustrated above is asynchronous, but if we want, it is quite easy to define a synchronous alternative using a technique we have already seen. A predicate \texttt{wait\_for\_answer/3} can be defined like so:

\begin{verbatim}
wait_for_answer(Pid, Answer, Options) :-
    receive({
        Answer if arg(1, Answer, Pid) ->
            true 
    }, Options).	
\end{verbatim}

\noindent Then \texttt{toplevel\_call\_synch/3-4} can be implemented as follows:

\begin{verbatim}
toplevel_call_synch(Pid, Goal, Answer) :-
    toplevel_call_synch(Pid, Goal, Answer, []).

toplevel_call_synch(Pid, Goal, Answer, Options) :-
    toplevel_call(Pid, Goal, Options),
    wait_for_answer(Pid, Answer, Options).
\end{verbatim}

\noindent and \texttt{toplevel\_next\_synch/2-3} like so:

\begin{verbatim}
toplevel_next_synch(Pid, Answer) :-
    toplevel_next_synch(Pid, Answer, []).

toplevel_next_synch(Pid, Answer, Options) :-
    toplevel_next(Pid, Options),
    wait_for_answer(Pid, Answer, Options).
\end{verbatim}

\noindent Note that since \texttt{toplevel\_stop/1} does not produce a response message, no \texttt{toplevel\_stop\_synch/2-3} is defined, and \texttt{toplevel\_stop/1} can be used as it is.

Here is a simple test that shows how it works:

\begin{verbatim}
?- toplevel_spawn(Pid).
Pid = 67590967.
?- toplevel_call_synch($Pid, mortal(Who), Answer, [
       template(Who),
       limit(1)
   ]).
Answer = success(67590967, [socrates], true).
?- toplevel_next_synch($Pid, Answer).
Answer = success(67590967, [plato], true).
?- toplevel_next_synch($Pid, Answer).
Answer = success(67590967, [aristotle], false).
?- 
\end{verbatim}

\subsubsection{Reconstructing findall/3 and findnsols/4}\label{sec:reconstructing-findall}

In order to specify the semantics of a new Prolog predicate API it can be a good idea to implement built-in Prolog predicates with a well-known semantics using the API. For example. if we did not already have a built-in predicate \texttt{findall/3}, using a toplevel actor we could have implemented it like so:

\begin{verbatim}
findall(Template, Goal, List) :-
    toplevel_spawn(Pid, [
        session(false)
    ]),
    toplevel_call(Pid, Goal, [
        template(Template)
    ]),
    receive({
        success(Pid, List, false) ->
            true ;
        failure(Pid) ->
            List = [] ;
        error(Pid, Error) ->
            throw(Error)
    }).
\end{verbatim}

\noindent It may not be useful, but it does say something about the relation between \texttt{findall/3} and the \texttt{toplevel\_*} predicates.\footnote{Here, we were inspired by \cite{DBLP:conf/cl/Tarau00}.} Of course, implementing these predicates in separate processes makes absolutely no sense. Examples such as these only exercises some of the capabilities of the \texttt{toplevel\_* API} and may seem somewhat pointless. Concurrency and distribution is where they really shine.

We can take this one step further by also implementing \texttt{findnsols/4}, which works like \texttt{findall/3}, but generates at most N solutions to the specified goal. %Thus, the length of the list of solutions will not be greater than N. 
This predicate is especially useful if the goal may have an infinite number of solutions.

\begin{verbatim}
findnsols(N, Template, Goal, List) :-
    toplevel_spawn(Pid, [
        session(false)
    ]),
    toplevel_call(Pid, Goal, [
        template(Template),
        limit(N)
    ]),
    wait_solutions(Pid, List).

wait_solutions(Pid, List) :-
    receive({
        success(Pid, List, false) ->
            true ;
        success(Pid, List0, true) ->
            (   List = List0 
            ;   toplevel_next(Pid),
                wait_solutions(Pid, List) 
            ) ;
        failure(Pid) ->
            List = [] ;
        error(Pid, Error) ->
            throw(Error)
    }).
\end{verbatim}

\noindent Ciao Prolog provides this predicate, but only in a deterministic version. The version provided by SWI-Prolog is nondeterministic, and so is our reconstruction. We borrow an example of its use from the SWI-Prolog manual:\footnote{\url{https://www.swi-prolog.org/pldoc/doc\_for?object=findnsols/4}}

\begin{verbatim}
?- findnsols(5, I, between(1, 12, I), L).
L = [1, 2, 3, 4, 5] ;
L = [6, 7, 8, 9, 10] ;
L = [11, 12].
?-
\end{verbatim}

\noindent Interestingly, it turns out that the version provided by SWI-Prolog serves as an important building block in our proof-of-concept implementation of the \texttt{toplevel\_*} predicates. (It was implemented in SWI-Prolog by Jan Wielemaker in order to support \texttt{library(pengines)}, the library on top of which SWISH is built.)



\section{Prolog nodes are also agents}\label{sec:nodes-as-agents}
As is evident from our agent taxonomy in Figure~\ref{fig:prolog-agent-taxonomy}, we regard a Prolog node as a kind of Prolog agent. We consider it as such because, as we shall see, a node incorporates a significant subset of the capabilities of a toplevel actor. In particular, it enables clients to submit queries and to receive solutions lazily -- one at a time (or, if required, more than one at a time), in the order that Prolog finds them.

The diagram in Figure~\ref{fig:a-node} illustrates the anatomy of an executing Prolog node, which is accessed by a user-client that knows its URI and is authorized to utilize its capabilities. 

\begin{figure}[h]

\centering
	\includegraphics[width=9cm]{a-node}
    \caption{The anatomy of a Prolog node. A terminal running in a web browser is attached to it.}
    \label{fig:a-node}
\end{figure}

\noindent As suggested by the diagram, a node is comprised of components such as a node controller, web APIs, a shared database containing predicates defined by the owner, a container serving as the ``home'' of actors such as echo actors, supervisors and toplevels, and various settings for the owner to configure the node. As we shall see, not all nodes include all these components.

\subsection{The node controller}\label{sec:actor-manager}

In the most capable profile, the node controller is entrusted with a multitude of responsibilities designed to ensure the smooth and secure operation of the node. Although not every node controller undertakes all of these tasks, the most advanced implementations are expected to manage web API connections, authenticate and authorize clients, and spawn actors upon client request. Additionally, they impose various resource restrictions to maintain optimal performance, oversee registries, monitors, and links, and guarantee that messages sent back to clients adhere to the required format. The node controller also terminates actors that have lost client connections and implements shell utilities such as the \$Var substitution mechanism and the redefinition of I/O predicates.

From an external perspective and when dealing with actors, the intricate workings of the node controller are intended to remain largely concealed, functioning much like an invisible ``ether'' through which messages seamlessly flow. In this manner, the node controller abstracts away the underlying technical details, allowing users to benefit from its comprehensive management capabilities without the need to engage directly with its multifaceted processes.



\subsection{The shared Prolog database}\label{sec:shared-db}

A node may host a shared Prolog database accessible over HTTP. It can be:

\begin{itemize}
  \item \emph{self-contained} - all predicates required for the exported relations are present locally, making the dataset directly consumable by other nodes or clients;
  \item \emph{dependent} - code or data rely on external modules or remote predicates, which reduces portability unless those dependencies are resolvable at the destination.
\end{itemize}

\noindent The predicates -- which will be referred to as the \textit{node-resident} predicates -- are typically maintained by the owner of the node. External users are not able to make any changes to them, only the owner of the node is allowed to do so. Any client connected with the node has access to these predicates in addition to the Web Prolog built-in predicates, and can pose queries over the web APIs formulated in terms of them. 

The contents of the file \texttt{mortal.pl} must be valid Web Prolog source code. It may be something like this:

\begin{verbatim}
mortal(Who) :- human(Who).

human(socrates).
human(plato).
human(aristotle).
\end{verbatim}

\noindent This is also what a user is presented with when steering an ordinary browser to the URI  \texttt{http://n7.org}. This is an example of a shared Prolog database that is  \textit{self-contained}. In Chapter~\ref{sec:the-prolog-web} we will demonstrate how this property makes it possible to achieve what we think of as \textit{instant portability}.

However, a file might also look like this:

\begin{verbatim}
:- use_module(philosophers, [philosopher/1]).

mortal(Who) :- human(Who).

human(socrates).
human(Who) :- philosopher(Who).
\end{verbatim}

\noindent Since the clauses for \texttt{philosopher/1} are not available, this database is \textit{dependent}.

\subsection{Deploying a node}\label{sec:node-setup}

%Deployment involves: (i) exposing one or more of the three web APIs; (ii) publishing a shared database under stable URIs; and (iii) selecting a \emph{profile} to advertise to clients. An implementation can implement more than it advertises; what matters to clients is what the running instance \emph{exposes}.
%
%---

\noindent To make this as concrete as possible, let us assume that the owner of a node wants to make the data in the file \texttt{mortal.pl} available to clients. The owner uses the following command to set up a Prolog node:

\begin{verbatim}
$ prolog-node --port=80 --src=mortal.pl --settings=settings.pl
\end{verbatim}

\noindent As a result, a new node is created and the contents of the file is loaded into its shared Prolog database. The node is configured according to the settings given in the file \texttt{settings.pl}. After having performed all the usual moves necessary for the deployment of a web server, access to the node is offered through the set of API endpoints listed in Table~\ref{table:node-urls}:

\renewcommand{\arraystretch}{1.2}
\begin{table}[h]
\begin{center}
    \begin{tabular}{ | p{5cm} | p{6cm} |}
    \hline
    \textbf{API endpoint} &  \textbf{Resource or functionality} \\ \hline
    \texttt{http://n7.org} & The shared Prolog database \\ \hline
    \texttt{http://n7.org/shell} & A simple shell environment\\ \hline    
    \texttt{http://n7.org/call} & The stateless HTTP API \\ \hline
    \texttt{http://n7.org/toplevel\_spawn} & The semi-stateful HTTP API \\ \hline
    \texttt{ws://n7.org/actor} & The stateful WebSocket API \\ \hline 
    \end{tabular}
\end{center}
\caption{Proposal for how to link URIs with targets. Note that the stateful and semi-stateful APIs have additional endpoints.}
\label{table:node-urls}
\end{table}

\noindent By steering an ordinary browser to the URI \texttt{http://n7.org} an authorized user may want to just browse the program. Or perhaps query it by making a visit to \texttt{http://n7.org/shell}. Or perhaps the plan is to build an end-user application talking to the node over any of its web APIs.

Note also that because a Prolog node is also a web server, it may offer other HTTP endpoints as well, some of which may not be relevant to its core functionality.


Other ecosystems show that this strategy works. In the Prolog world, Jan Wielemaker's \textit{Cliopatria} demonstrated how a substantial Web application could be delivered as a ready-to-run system: users provided their data, made modest extensions, and obtained a full Semantic Web portal. In the Elixir world, Phoenix goes even further by generating complete, fault-tolerant applications with supervision, telemetry, routing, and deployment already in place. Developers simply add the behaviour that makes the system theirs. These examples illustrate how a language ecosystem can supply standardised, reproducible scaffolds that capture the best practices of a community.


\subsection{Settings}\label{sec:node-settings}

As suggested by Figure~\ref{fig:a-node}, a Web Prolog node is governed by a set of \emph{owner-controlled settings}. Only the node owner may change these settings; some of them may be \emph{readable} by clients over HTTP, so that clients can discover the node's advertised capabilities and operational limits. Collectively, the settings determine which capabilities the node offers, who may access them, and the resource envelope within which actors and sessions are allowed to run.

It is useful to group settings into five categories:

\begin{enumerate}
\item \textbf{Identity and advertised capability.} The node advertises a \textit{profile} to clients. This is primarily a discovery and compatibility mechanism: it tells clients what kind of service they can expect.

\item \textbf{Resource governance.} The node may enforce limits on injected code size, actor-local database size, cache bounds, and the maximum number of concurrent actors that client code is allowed to spawn. At the transport layer it may also enforce connection-related limits (for example, the maximum number of connections, and the maximum number of agents sharing a connection). These limits are part of the node's resource policy and are set by the node owner.

\item \textbf{Admission control.} The node may restrict which clients may connect and which goals may be executed, for instance via allow/deny lists and related authorization rules.

\item \textbf{Session lifecycle and PTCP timeouts.} Two timeouts are particularly important for PTCP sessions (cf.\ the PTCP state diagram in Figure~\ref{fig:statechart-2}). A \emph{running timeout} bounds how long a session may remain in the working state (i.e \textbf{s2}). If this timeout is exceeded, the node aborts the current goal and returns an error (e.g.\ \texttt{timeout\_exceeded}) to the caller, while keeping the session actor alive. An \emph{idle timeout} bounds how long a session may remain inactive in an idle state (i.e \textbf{s1} or \textbf{s3}). If this timeout is exceeded, the node terminates the session actor. If the session is monitored, the client will subsequently observe a \texttt{down} message consistent with normal actor termination.

\item \textbf{Bundle loading policy.} The node may require, or merely prefer, that loaded bundles are self-contained, i.e.\ do not rely on unresolved external dependencies.
\end{enumerate}

\noindent Note that readable over HTTP does not mean modifiable: the intention is that the node may expose a read-only subset of the settings that are relevant to interoperability and capacity planning, while keeping sensitive policy (notably admission rules) private.

Several of these settings are best understood as \emph{capability gates} rather than mere tuning parameters. In particular, a node distinguishes between what the owner may do and what external clients may do. Clients are intended to program \emph{session-local} concurrent computations, while the owner configures the long-lived surface of the node and installs durable services. This distinction will later be reflected in which concurrency-oriented predicates and options are admitted for client code, especially for lifetime control (for example whether a spawned actor may outlive its parent) and, in the distributed case, which combinations of remote placement and lifetime options are permitted.

\subsection{Two transports -- three web APIs}\label{sec:transports-apis}

The primary responsibility of the node controller is to facilitate web API interactions. The most capable nodes are designed to support two different transports -- HTTP and WebSocket -- and three distinct API types: a stateless HTTP API, a semi-stateful HTTP API, and a stateful WebSocket API. This architecture enhances flexibility by increasing the likelihood that at least one API will effectively serve a given end-user application.\footnote{Strictly speaking, HTTP and WebSockets both live in the application layer, not the transport layer. But in practice, even if TCP is the real transport protocol under the hood, it makes sense to think of them as alternative ways of transporting data between agents.}


\subsubsection{The HTTP transport}\label{sec:http-transport}

The \textit{Hypertext Transfer Protocol} (HTTP) is a core communication protocol of the World Wide Web, enabling the exchange of resources like web pages and data between clients and servers. Operating as a stateless request--response protocol, it supports methods such as GET, used to retrieve resources, and POST, used to submit data to a server. 

A client interacting with a node by making requests over a HTTP connection is illustrated in Figure~\ref{fig:http-communication}.

\begin{figure}[h]
    \centering
	\includegraphics[width=8cm]{http-communication}
    \caption{Interaction between a client and a node over an HTTP connection.}
    \label{fig:http-communication}
\end{figure}


\noindent HTTP typically uses TCP for reliable transmission and supports secure communication through HTTPS with TLS encryption. Evolving since the 1990s, HTTP has advanced through versions like HTTP/2 and HTTP/3, enhancing performance and features like caching, authentication, and content negotiation. There is a vast infrastructure for HTTP and HTTPS that already exists (proxies, firewalls, caches, and other intermediaries), ensuring that security requirements of the modern web are fulfilled and the browser and the node can validate each other.

%\todo{Point to examples in Appendix.}


\subsubsection{The WebSocket transport}\label{sec:websocket-transport}

In order to enable a client to control all aspects of a set of toplevels or other kinds of actors, the most capable type of node offers a WebSocket sub-protocol. WebSocket is a real-time, low latency, bidirectional protocol for asynchronous communication between a client and a node. A WebSocket sub-protocol is an application-layer protocol, which operates on top of the WebSocket protocol, leveraging the transport services provided by TCP. The three phases in the life-cycle of a WebSocket connection is illustrated in Figure~\ref{fig:ws-communication}.

\begin{figure}[h]
    \centering
	\includegraphics[width=8cm]{ws-communication}
    \caption{The life-cycle of an interaction between a client and a node over a WebSocket connection.}
    \label{fig:ws-communication}
\end{figure}

In the first phase (in blue in the diagram) a WebSocket connection is initiated by the client via an HTTP connection which is then switched over to the WebSocket protocol. In the second phase (in black), the client and the node starts talking to each other. This is where the sub-protocol comes in. 

Subject to the value of a HTTP header or parameter, answers and other kinds of messages arriving from the node are returned in the form of either JSON or Prolog text. Client code written in languages other than Prolog would normally request that they be encoded in JSON. 

WebSocket client libraries are available for the most common general programming languages. However, the most common use of WebSockets is in communication between a web browser and a server. In all the major web browsers, the WebSocket object provides a JavaScript API for creating and managing a WebSocket connection to a server, as well as for sending and receiving data on the connection.\footnote{See \url{https://www.w3.org/TR/websockets/}} The role of JavaScript here is to create a new websocket, to define the necessary handlers for \texttt{onopen}, \texttt{onmessage}, \texttt{onerror} and \texttt{onclose} messages, and to call the methods \texttt{send} or \texttt{close} for sending messages or closing the connection.%A synopsis that leaves out a lot of details can be given as follows:
\footnote{For a good general introduction to WebSocket we recommend \cite{lombardi2015websocket}.}

\medskip
\noindent In the next section, and with the help of a number of scenarios, the actual APIs will be demonstrated. For details, see Appendix~\ref{sec:manual}.

\section{Node profiles}\label{sec:node-profiles}

Not all Prolog nodes are equal. They vary in the capabilities they offer
clients, and we use the term \emph{profile} -- borrowed in the standards
sense -- for the declared set of capabilities a node commits to providing.
A profile is a tailored subset of the full design: it constrains what a node
must offer in order to make that offer verifiable, portable, and predictable.
Less capable profiles are easier to specify, implement, and standardise,
while still participating in the same ecosystem.\footnote{Throughout this
section \emph{capability} means a client-visible feature or power provided
by a node -- support for a given web API or built-in predicate. The narrower capability-security sense of an unforgeable authority token
-- a value whose possession confers the ability to act -- is developed in
Chapter~\ref{sec:the-prolog-web}. The two uses are related: node-level
capabilities determine what actions are available in principle, while
capability values (such as pids) determine which of those actions a client
can actually perform.}

We distinguish two kinds of capability: \emph{language capabilities} (which
Prolog fragments and predicates are available) and \emph{interactional
capabilities} (which transports and web APIs the node supports). Four
profiles are defined: RELATION, ISOBASE, ISOTOPE, and ACTOR. Their
nesting is shown in Figure~\ref{fig:node-profiles}: features in the inner
rings are required by every profile, while features in the outer rings are
restricted to more capable ones.

\begin{figure}[h]
    \centering
    \includegraphics[width=10cm]{profile-hierarchy-new-2}
    \vspace{0.3cm}
    \caption{The profile hierarchy. Each ring adds capabilities to those of
             the profiles inside it.}
    \label{fig:node-profiles}
\end{figure}

The RELATION profile is the base. It requires only the stateless HTTP API
and a set of owner-defined relations; it does not require any Web Prolog
built-in predicates. ISOBASE adds the standard built-in predicate set,
\texttt{rpc/2-3}, and asynchronous RPC. ISOTOPE further adds a
session-oriented toplevel actor accessible over the semi-stateful HTTP
API, enabling persistent state and interactive I/O within a session. ACTOR
adds the full Erlang-style concurrency predicates -- \texttt{spawn/2-3},
\texttt{!/2}, \texttt{receive/1-2}, and related primitives -- and exposes
them over a bidirectional WebSocket connection. A node offering the ACTOR
or ISOTOPE profile can be both queried and programmed; a node offering only
ISOBASE or RELATION can only be queried.

The RELATION profile differs from the others in kind, not merely in degree.
An ISOBASE, ISOTOPE, or ACTOR node accepts arbitrary Prolog goals -- within
the predicate subset defined by the profile -- and evaluates them. A
RELATION node does not. It exposes specific named relations through a fixed
query interface: a client can ask \texttt{wife(H,\,W)} because the node
advertises that relation, not because it evaluates arbitrary Prolog code.
The goal in a request to a RELATION node is better understood as a pattern
matched against a known schema than as a program to execute. A conformant
RELATION node need not contain a Prolog engine at all.

A node is \emph{conformant} at a given profile if it provides authorised
clients with every capability that profile specifies. A node \emph{may}
provide more -- additional predicates, extra web APIs -- but is not
required to. A RELATION node may, for example, support \texttt{rpc/2-3}
and the semi-stateful HTTP API, but only the stateless API is mandated.
Clients that depend on extras sacrifice portability across nodes at that
profile level.

The hierarchy is intentionally non-exhaustive. One can also imagine
\emph{orthogonal} profiles that cut across the capability ladder by adding
restrictions or guarantees rather than extending it -- a PURE profile
limiting clients to declarative Prolog, for instance, or a statechart
discipline as an ACTOR sub-profile. Section~\ref{sec:other-profiles}
returns to this possibility.


\subsubsection*{What are profiles good for?}\label{sec:profiles-why}

Profiles do two things simultaneously: they give clients a stable contract
and give node implementers a scalable commitment. A client talking to an
ACTOR node knows it may spawn actors, use persistent sessions, and receive
push events over WebSocket; the same client talking to an ISOBASE node
knows it may not. A node implementer choosing to deploy at the RELATION
level knows that a Python microservice wrapping a database is sufficient;
one targeting ACTOR knows what the full engineering cost entails. The
profile a node advertises is thus a public, verifiable statement of what
both parties can assume.

Crucially, this contract operates at the level of \emph{observable
interaction}, not implementation. Across all four profiles, the same
shell-style pattern is preserved: a client submits a goal and receives
solutions one at a time, lazily, on demand. What differs is only the
infrastructure behind that pattern -- stateless HTTP, a semi-stateful
session actor, or a live WebSocket stream. In the extreme case a
RELATION node may produce correct answers without any Prolog execution at
all; from the client's perspective the interaction is indistinguishable.
This separation of observable contract from implementation commitment is
what allows the profile hierarchy to be both strict (conformance is
testable) and open (any conformant implementation qualifies, regardless
of how it works internally).

Profiles also carry a security dimension, developed fully in
Section~\ref{sec:security-prolog-web}. A node's profile determines not
only what a client can \emph{do} but what it can \emph{express}: excluding
I/O and dynamic database modification from ISOBASE is not merely a
protocol convenience but removes those operations from the
client-accessible language fragment entirely. Greater expressiveness at
higher profiles therefore comes paired with greater exposure, a tradeoff
the security architecture addresses through inexpressibility, name
discipline, and lifetime containment. The profile a node advertises is,
in this sense, also a declaration of its security posture.

The following subsections work through five scenarios -- RELATION, ISOBASE,
ISOTOPE, ACTOR via browser, and ACTOR via social robot -- each illustrating
the interaction model that its profile enables.


\section{Node profiles -- OLD}\label{sec:node-profiles}

Not all Prolog nodes are equal; they vary in the capabilities they offer clients. We call the set of capabilities that a node provides its \emph{profile}. We borrow the term \emph{profile} in the standards sense: a tailored subset that constrains and organises a fuller design to meet particular goals (simplicity, performance, security, portability). Less capable profiles are easier to specify, implement, and standardise -- while still participating in the same ecosystem.\footnote{Throughout this
section \emph{capability} means a client-visible feature or power provided
by a node -- support for a given web API or built-in predicate. The narrower capability-security sense of an unforgeable authority token
-- a value whose possession confers the ability to act -- is developed in
Chapter~\ref{sec:the-prolog-web}. The two uses are related: node-level
capabilities determine what actions are available in principle, while
capability values (such as pids) determine which of those actions a client
can actually perform.}


In this chapter we distinguish between two broad kinds of capabilities: \textit{language capabilities} (what Prolog fragments and predicates are available) and \textit{interactional capabilities} (what transports and web APIs the node supports). Four profiles will be presented and discussed in some detail: RELATION, ISOBASE, ISOTOPE and ACTOR. The Venn diagram in Figure~\ref{fig:node-profiles} provides an overview of how these profiles are related in our proposed design of the Prolog Trinity ecosystem. 

\begin{figure}[h]
    \centering
	\includegraphics[width=10cm]{profile-hierarchy-new-2}
    \vspace{0.3cm}\caption{A hierarchy of node profiles and the capabilities each level adds.}
    \label{fig:node-profiles}
\end{figure}

%\noindent Intuitively, ACTOR $\supset$ ISOTOPE $\supset$ ISOBASE $\supset$ RELATION: every ACTOR node is also an ISOTOPE node, and so on, but not vice versa. The RELATION profile is the least capable of the four.


For example, owner-defined shared predicates and the HTTP API are located inside the area marked RELATION to signify that these features are supported by \emph{all} profiles, whereas placing the WebSocket API in the outermost area indicates that it is only available in the ACTOR profile.

In short, when running against a node that offers the ACTOR or ISOTOPE profile, we can both query and program it. We may load predicates into a toplevel at the time of its creation, or use \texttt{assert/1} and \texttt{retract/1} to update its private database at any time during its lifecycle. Only the ACTOR profile, however, allows us to use predicates such as \texttt{spawn/2-3}, \texttt{!/2} and \texttt{receive/1-2} in queries and programs. Running against an ISOBASE or RELATION node, we can only query it, and of these two only the ISOBASE profile allows us to pose queries using built-in predicates. 

The RELATION profile differs from the other three in kind, not merely in degree. An ISOBASE, ISOTOPE, or ACTOR node accepts arbitrary Prolog goals -- within the predicate subset defined by its profile -- and evaluates them. A RELATION node does not. It exposes specific named relations through a fixed query interface: the client can ask \texttt{wife(H,\,W)} because the node advertises that relation, not because it evaluates arbitrary code. The goal parameter in a request to a RELATION node is better understood as a query pattern against a known schema than as a program to execute. This is what makes it possible for a RELATION node to be implemented without a Prolog engine at all.

To be considered \textit{conformant}, a node \textit{must} provide authorized clients with \textit{all} capabilities specified by its profile. For instance, an ACTOR node \textit{must} offer clients access to the three web APIs and the entire set of built-in predicates indicated in Figure~\ref{fig:node-profiles}. Similarly, an ISOTOPE node \textit{must} provide clients with the two HTTP APIs and all predicates defined in the ISOTOPE subset of Web Prolog. Other profiles follow analogous requirements.

A node \textit{may} provide \textit{more} than what is mandated by its profile, including additional predicates or web APIs.  For example, a RELATION node may support the \texttt{rpc/2-3} predicate as well as both the stateless and semi-stateful HTTP APIs, but it is not obligated to do so. Only the stateless API is required. A dependency on extra features comes with a price: it may make instant portability impossible. [TODO: Elaborate this - but perhaps not only here]

The profile hierarchy presented here is intentionally non-exhaustive. Beyond adding new rungs to the capability ladder, one can also imagine \emph{orthogonal} profiles that cut across the hierarchy by adding restrictions or guarantees. Section~\ref{sec:other-profiles} briefly discusses this possibility.

%The profile hierarchy presented here is not intended to be exhaustive; Section~\ref{sec:other-profiles} discusses how it might be extended.

%The profile hierarchy presented here is not intended to be exhaustive. One could, for example, imagine a further profile --- perhaps \textsc{PURE} --- that restricts clients to a fully declarative subset of Prolog, excluding constructs such as cut and non-pure arithmetic. Such a profile would be orthogonal to ISOBASE and ISOTOPE, and could still admit features like source code injection. We do not pursue this direction further here; the purpose of mentioning it is merely to illustrate that profiles are a flexible structuring device rather than a closed taxonomy.


\subsubsection*{What are profiles good for?}\label{sec:profiles-why}

The purpose of node profiles is to make explicit which interactional guarantees a client may rely upon, and which obligations a node commits to satisfy. An ACTOR node, for example, requires more of the toplevel actor at its core than an ISOTOPE node does: it must support full actor-style concurrency and bidirectional interaction, whereas an ISOTOPE node need only support a single session-oriented toplevel governed by PTCP and acting as a shell.

Across all profiles, however, a common interaction pattern is preserved. In each of the scenarios presented above, the shell allows the user to submit a query and to lazily backtrack through solutions, one at a time. What differs between profiles is not this external behavior, but how much infrastructure is required to support it. The same shell-style interaction can be realized over stateless HTTP, semi-stateful HTTP, or a stateful WebSocket connection, typically involving one network roundtrip per solution. Profiles therefore separate \emph{observable interaction} from \emph{implementation commitment}. In the extreme case, a RELATION node may have no Prolog implementation at all -- the observable interaction is identical, but the implementation commitment is merely to produce correct answers, not to provide a Prolog execution environment.

This distinction matters because the Prolog shell's interaction model is central to logic programming. Presenting one solution at a time, pausing for user input, and resuming search on demand exposes the mechanics of backtracking and makes nondeterministic computation tractable and inspectable. Profiles ensure that this familiar interaction can be offered even by nodes with limited capabilities, while allowing more capable nodes to support richer forms of control, state, and concurrency.

In this way, profiles serve two complementary goals. For clients, they provide a stable contract describing what can be assumed about interaction, state, and control. For node implementers, they offer a way to scale complexity: simpler profiles are easier to implement and deploy, while more powerful profiles enable advanced applications without forcing those costs on every node in the ecosystem.

Profiles also have a security dimension that becomes fully explicit only in
Chapter~\ref{sec:the-prolog-web}. A node's profile determines not only what a
client can \emph{do} but what it can \emph{express}: the ISOBASE profile excludes
I/O and dynamic database modification not merely because these are unnecessary
for stateless interaction, but because excluding them removes entire classes of
potentially dangerous operations from the client-accessible language fragment.
The ACTOR profile's greater expressiveness therefore comes paired with greater
exposure -- a tradeoff that the security architecture of Section~\ref{sec:security-prolog-web}
addresses through inexpressibility, name discipline, and lifetime containment.
The profile a node advertises is thus also a public statement of its security posture:
it tells clients and node owners alike where the capability boundary lies.

%\subsubsection*{What are profiles good for?}\label{sec:profiles-why}
%
%An ACTOR node requires more of the toplevel actor running at its core than is required of the toplevel at the core of an ISOTOPE node. 
%
%---
%
%As we have seen, in each of these scenarios the shell allows the user to submit queries and lazily backtrack through solutions using the most powerful of the web APIs that the node supporting this profile admits. (Important!)
%
%It has been demonstrated that this form of interaction can be held over either stateless HTTP, semi-stateful HTTP or stateful WebSocket, making one network roundtrip for each solution.
%
%---
%
%The Prolog shell's interaction model is a core feature of logic programming, revealing Prolog's internal workings and its reliance on backtracking. This model helps users refine problem-solving, write efficient code, and navigate complex searches while building an intuitive grasp of logical inference.
%
%This interactive approach closely ties to backtracking, where Prolog revisits decision points, retracts choices, and explores alternative paths. The shell's step-by-step behavior, offering one solution at a time and pausing for input, exemplifies this process. When requesting another solution, Prolog deliberately fails the current path, triggering backtracking to find new branches. This iterative process highlights Prolog's method of pruning and branching the solution space.
%
%This model is both educational and practical. It helps users visualize core logic programming concepts, like deriving multiple answers from one query, and provides tools for refining programs. By observing solutions, users adjust rules, re-run queries, and align logic with desired outcomes. Immediate feedback enables continuous learning and improvement.
%
%Understanding the Prolog shell and its link to backtracking enhances problem-solving, strengthens conceptual knowledge, and supports the creation of efficient solutions.

%\section{Node profiles in scenarios}\label{sec:node-scenarios}

The following subsections present several scenarios illustrating interactions between end-user clients, nodes, and actors. The scenarios aim to demonstrate the various types of communication that can occur, starting with the loading of a Prolog shell -- a specialized terminal implemented in JavaScript and executed within a web browser.

\subsection{Browser talking to a RELATION node}\label{sec:browser-to-relation-node}

As indicated by the profile hierarchy in Figure~\ref{fig:node-profiles}, the stateless HTTP API is the only web API that \textit{must} be supported by a RELATION node. This does not distinguish it from nodes having a more capable profile, as this is required by them too. What makes it stand out is that it is not required to offer access to any Web Prolog built-in predicates or control constructs \textit{at all}. Even control constructs such as \texttt{,/2} and \texttt{;/2}, may not be available. 

In the scenario depicted in Figure~\ref{fig:browser-to-relation-node}, a user equipped with an ordinary web browser is talking to the RELATION node at \texttt{http://n1.org}. The web-based terminal displays the remnants of a session during which the user executed two queries. The two question marks signify that neither the owner-defined resources nor the way answers are produced is known by the user.

\begin{figure}[h]
    \centering
	\includegraphics[width=9cm]{browser-to-relation-node}
    \vspace{0.2cm}\caption{A scenario involving a web browser in conversation over HTTP with a remote RELATION node. The two question marks signify that neither the owner-defined resources nor the way answers are produced is known by the user.}
    \label{fig:browser-to-relation-node}
\end{figure}

\noindent The interaction started with the loading of the Prolog shell at \texttt{http://n1.org/shell}. 
%This is when the welcome banner was displayed, along with the \texttt{?-} prompt inviting the user to a session. A JavaScript process started running a loop tasked with reading queries from the terminal, transforming them into stateless HTTP requests, picking up responses, and writing answers to the terminal. 
The interaction opened with the welcome banner and a \texttt{?-} prompt; a JavaScript loop then began reading queries, translating them into HTTP requests, and writing answers back to the terminal.

Let us look at how the first solution to \texttt{?-wife(H,W)} was retrieved by using the stateless HTTP API. The URI pointing to this solution is:

\begin{verbatim}
http://n1.org/call?goal=wife(H,W)&limit=1&offset=0
\end{verbatim}

\noindent Note the use of \texttt{limit} and \texttt{offset} parameters. Although here they are parameters rather than options, their defaults and the effect that they have on answer terms is the same as for the options accepted by \texttt{toplevel\_call/3}. 

URIs such as these are simple, they are meaningful, they are declarative, they can be bookmarked, and responses are cacheable by intermediates. Most of these desirable properties are there because the client-node interaction involved during the process of resolving the URIs is stateless. 

When the user entered the query \texttt{?-wife(H,W)}, (and assuming that a relative path was used instead of an absolute one) the following GET request was made to the node:

\begin{verbatim}
GET /call?goal=wife(H,W)&limit=1&offset=0 HTTP/1.1
Accept: application/json
\end{verbatim}

\noindent The \texttt{Accept} header requested responses in JSON. (In this case the header could have been left out since JSON is the default.)\footnote{A \texttt{format} parameter is also supported as this allows us to easily make a request from the address field of a web browser.}

The node responded with the following JSON structure, indicating that the query succeeded: 

\begin{verbatim}
{ "type": "success",
  "data": [{"H":"socrates","W":"xantippa"}], 
  "more": true }
\end{verbatim}

\noindent The \texttt{data} field contains the variable bindings and the value \texttt{true} of the \texttt{more} field signals that further solutions may be available. This triggered the printing of the variable bindings and served as a prompt for how to continue the exploration.

The user wanted to know if there were more solutions and typed a semicolon which led to the following request being made, looking exactly like the previous request, except that the \texttt{offset} parameter was incremented from \texttt{0} to \texttt{1}:

\begin{verbatim}
GET /call?goal=wife(H,W)&limit=1&offset=1 HTTP/1.1
\end{verbatim}

\noindent The node responded like so:

\begin{verbatim}
{ "type": "success",
  "data": [{"H":"aristotle","W":"pythias"}], 
  "more": false }
\end{verbatim}

\noindent The \texttt{more} field was now \texttt{false}, indicating that no further solutions existed, and after having printed the next set of bindings and a full stop, the shell again produced the \texttt{?-} prompt, inviting the user to enter another query.

At this stage of the interaction, it is reasonable for the user to expect that \texttt{wife/2} has the mode \texttt{wife(?,\,?)} and that queries such as \texttt{?-wife(socrates,\,W)} or \texttt{?-wife(H,\,pythias)} should also succeed. A query such as \texttt{?-wife(plato,\,W)} should fail, however, and the response would be the following:

\begin{verbatim}
{ "type": "failure" }
\end{verbatim}

\noindent Finally, an erroneous query, such the attempt in the scenario to query a predicate which is not defined, should raise an error, and show up as a response like so:

\begin{verbatim}
{ "type": "error",
  "data": "Unknown procedure: =/2" }
\end{verbatim}

\noindent We can summarize this by saying that a RELATION node must offer clients a truly relational query interface over stateless HTTP. Any node failing to meet this expectation is not a conformant RELATION node. A RELATION node is required to respond to a client as if it were Prolog, thereby allowing the user to lazily backtrack over possible solutions. Although a RELATION node may not truly perform backtracking, it must behave externally as though it does. 

\medskip
\noindent The lack of built-in predicates naturally makes a RELATION node very easy to implement. There are two quite different ways a RELATION node can come into existence. First, a node that already implements a richer profile can be turned into a RELATION node by installing a filter that inspects the goal parameter in each \texttt{/call} request and only allows it through if it matches a declared relation. Such a filter can be as simple as matching the submitted goal against a whitelist of terms -- for instance, accepting only \texttt{employee(\_,\,\_,\,\_)} and \texttt{name(\_,\,\_)}, and rejecting everything else with an error. A node that is supposed to serve a particular application rather than act as a general compute server will often operate this way: it implements ISOBASE or more internally, but presents a RELATION interface to its clients.

Second, a RELATION node can be implemented from scratch in any programming language. (Recall that we do not require that a Prolog node is written in Prolog, only that it can talk Prolog.) At a minimum, an implementation must parse Prolog terms and produce variable bindings for queries against its data. A Python web service backed by a relational database, or a microservice wrapping a key--value store, can present itself as a conformant RELATION node as long as it implements the stateless HTTP API and responds with the expected JSON structure. This is the profile that makes the Prolog Web interoperable with the rest of the Web: any structured data endpoint can, in principle, participate in the ecosystem. 

%The lack of built-in predicates naturally makes a RELATION node very easy to implement. A key motivation for allowing a RELATION node to expose only such a minimal functionality is to make a Prolog Web node easy to implement in virtually any programming language. (Recall that we do not require that a Prolog agent is written in Prolog, only that it can talk Prolog.) At a minimum, an implementation must parse Prolog terms and create variable bindings for a given query. 


\subsection{Browser talking to an ISOBASE node}\label{sec:example-isobase}

For a node to adhere to the requirements of the ISOBASE profile, it \textit{must} support the same stateless HTTP API as required by the RELATION profile. In our next scenario, illustrated in Figure~\ref{fig:browser-to-n1}, the user is running the same query as in the previous scenario. It means that we do not need to describe the details of requests and responses again. Instead, we focus on the language capabilities of the node and the contents of the shared Prolog database offered by it, and visible in the figure.  

\begin{figure*}[h]
    \centering
	\includegraphics[width=\textwidth]{browser-to-isobase-node}
    \vspace{0.3cm}\caption{Stateless interaction in a scenario involving a web browser and a remote node. }
    \label{fig:browser-to-n1}
\end{figure*}

\noindent The main difference compared to the previous scenario is that an ISOBASE node \textit{must} make a fairly large number of Web Prolog built-in predicates available, including a large subset of ISO Prolog predicates (recall the cloud of ISO Prolog predicates in Chapter~\ref{sec:web-prolog}), standard library predicates, definite clause grammars, as well as predicates for working with the Web. (The scenario differs from the previous one in that it allows the user to browse the source code defining the predicates offered by the owner of the node.)

An ISOBASE node must also support \texttt{rpc/2-3}, an easy-to-use predicate for making \textit{nondeterministic remote procedure calls} from one node to another over stateless HTTP, communicating in a way that is essentially \textit{synchronous}. In the scenario in Figure~\ref{fig:browser-to-n1}, the URI \texttt{http://n1.org} in the first argument of the call to \texttt{rpc/2} shows that it accesses the relation \texttt{wife/2} made available by the node \texttt{n1.org} in the previous scenario. 

In addition to the synchronous rpc/2-3, an ISOBASE node must also support \textit{asynchronous} RPC using two built-in predicates; \texttt{promise/3-4} and \texttt{yield/1-2}. It turns out that these three predicates can all be implemented on top of the stateless HTTP API. They will be dealt with in more detail in Chapter~\ref{sec:the-prolog-web}.

% Of course, the user could have made the query like so instead:

% \begin{verbatim}
% ?- rpc('http://n1.org', wife(H, W)).
% H = socrates, W = xantippa ;
% H = aristotle, W = pythias.
% ?-
% \end{verbatim}

% \noindent An ISOBASE node must also support asynchronous RPC using \texttt{promise/3-4} and \texttt{yield/1-2}:

% \begin{verbatim}
% ?- promise('http://n1.org', wife(H, W), Ref).
% Ref = 30308811.
% ?- yield($Ref, Answer).
% Answer = success([wife(socrates,xantippa), 
%                   wife(aristotle,pythias)], false).
% ?-
% \end{verbatim}

%Stateless interaction has one significant constraint: once a query is submitted, it cannot be aborted. In fact, no outside party, including the client who originally submitted the request, can communicate in any way with the remote computation once it has begun evaluation. A Web Prolog query may take a long time to execute, or may never return with an answer, and yet, once a query has been entered, it cannot be aborted by the user. Because stateless calls cannot be interrupted, ISOBASE nodes must enforce a running timeout. Instead, queries must always be executed under a time limit so that, typically after just a few seconds or even less, a client would be met with an error message. 

Stateless interaction has one significant constraint: once a query is submitted, neither the client nor any other party can abort it. Because there is no control channel back to the running computation, ISOBASE nodes must enforce a timeout.
This is what happened with the user's second query in the Figure~\ref{fig:browser-to-n1} scenario: it was met with a \texttt{Timeout\,exceeded} message. Browsing the contents of the node's shared database, the problem seems to be that the node owner forgot to add a DCG rule \texttt{s\,{-}->\,[]} that would have allowed the call to terminate.

Another limitation is that a truly stateless request-response protocol simply cannot support updates to the dynamic private database of an actor, nor I/O. It means that ISO predicates for modifying the dynamic database (e.g. \texttt{assert/1} and \texttt{retract/1}), or I/O (e.g. \texttt{read/1} and \texttt{write/1}), are \textit{not} available. %As we shall see, it is the nature of the protocol that constrains what \textit{can} be made available. 
This is what happened with the user's third query, which was met with an \texttt{Unknown\,procedure} error message when trying to assert the clause \texttt{p(a)}. The predicate \texttt{assert/1} is simply not available to a client talking to an ISOBASE node.



\subsection{Browser talking to an ISOTOPE node}\label{sec:browser-talking-to-isotop-node}

An ISOTOPE node offers clients a shell backed by a toplevel actor. (Importantly, this does not mean that the internal \texttt{toplevel\_*} predicates are directly exposed to programmers.) By using the process identifier of this actor as a session key in otherwise stateless HTTP requests, the node provides a semi-stateful API: each request is stateless in itself, but thanks to the pid they address the same long-lived toplevel actor.

The client's initial request can populate the actor's private database, which can then be updated dynamically throughout the session. Furthermore, although the client must initiate all interactions, the toplevel actor can not only respond to queries but can also prompt the client for further input or produce output at arbitrary moments. As long as the client keeps supplying the session identifier, this will work just fine.

%In the scenario in Figure~\ref{fig:browser-to-isotope-node} the client has established communication with a toplevel actor running on an ISOTOPE node, which is a node that must provide all the capabilities that an ISOBASE node has, and then some.

In the scenario in Figure~\ref{fig:browser-to-isotope-node}, the client established communication with a toplevel actor on an ISOTOPE node, which is a node that must provide all the capabilities that an ISOBASE node has, and then some.

\begin{figure}[h]
    \centering
	\includegraphics[width=\textwidth]{browser-to-isotope-node}
    \vspace{0.1cm}\caption{A scenario involving a web browser in conversation over semi-stateful HTTP with a remote ISOTOPE node running a toplevel actor.}
    \label{fig:browser-to-isotope-node}
\end{figure}

The request below instructed the node controller to spawn a toplevel actor running as a session, and to populate its private database with the value of the \texttt{load\_text} option. (Let us assume that the clause \texttt{q(X)\,:-\,p(X)} was picked up from an editor window that is not shown in the diagram.)

\begin{verbatim}
POST /toplevel_spawn HTTP/1.1
Content-Type: application/json

{ "options": "[session(true), load_text('q(X):-p(X).')]" }
\end{verbatim}

\noindent The node responded with a JSON structure containing the pid for the spawned toplevel actor:

\begin{verbatim}
{ "type": "spawned",
  "pid" : 56120932 }
\end{verbatim}

\noindent This produced the prompt \texttt{?-} in the terminal, inviting the user to enter a query. The query \texttt{?-q(X)} caused the JavaScript shell to make the following request to the node:

\begin{verbatim}
GET /toplevel_call?pid=56120932&goal=q(X)&limit=1 HTTP/1.1
\end{verbatim}

\noindent (Note that the pid must be sent along when querying the toplevel.) The request told the node controller to make a call to \texttt{toplevel\_call/3} and to convert the resulting answer term into a JSON structure that encodes the first solution:

\begin{verbatim}
{ "type": "success",
  "pid" : 56120932,
  "data": [{"X":"a"}],
  "more": true }
\end{verbatim}

\noindent Except for the pid, the JSON structure representing this solution is identical to the response when making a stateless request.

The value \texttt{true} of the \texttt{more} property signifies that more solutions may be available. Asking for the next solution results in a similar message, except that the value of the \texttt{data} property becomes \texttt{[\{"X":"b"\}]} and the value of \texttt{more} becomes \texttt{false}.


When the user asked for another solution, the following request was made:

\begin{verbatim}
GET /toplevel_next?pid=56120932 HTTP/1.1
\end{verbatim}

\noindent Note that it took the \textit{combination} of the clause that was loaded into the toplevel's private Prolog database, and data in the shared database to answer the query \texttt{?-q(X)}.

Because the toplevel persists across requests, it makes sense to allow the client to update its private database using \texttt{assert/1} and \texttt{retract/1} during the session.

\begin{verbatim}
GET /toplevel_call?pid=56120932&goal=assert(r(c,d)) HTTP/1.1
\end{verbatim}

\noindent With this, the clause \texttt{r(c,d)} ended up in the private database of \texttt{56120932}, but was removed again when the user decided to call \texttt{retract(r(X,d))}.

The user then called \texttt{greetings/0}, which calls \texttt{sleep(1)} and \texttt{writeln(hello)} in a loop. The \texttt{writeln/1} predicate has been redefined by the node controller as \texttt{writeln(Term)\,:-\,output(Term)} so this resulted in a series of messages of the following form to be sent back to the client, one per second:

\begin{verbatim}
{ "type": "output",
  "pid" : 56120932,
  "data": "hello" }
\end{verbatim}

\noindent For this to work over HTTP, a trick somewhat similar to so-called \textit{long-polling} is employed: every time a JSON structure of type \texttt{output} reaches the client, it must issue a new request for a response, which can come as more output, or some other response. See Section~\ref{sec:semi-stateful-notes} for more details.

The call to \texttt{greetings/0} will run forever unless the user does something about it. When the user hits Control-C, the nonterminating query is aborted, and the node responds with the following message:

\begin{verbatim}
{ "type": "abort",
  "pid" : 56120932 }
\end{verbatim}

\noindent If the client does not exit the remote process, the node may terminate it after a set time with no activity. The pid will no longer be valid and any changes made to the workspace of the process will be lost. 


\subsection{Browser talking to an ACTOR node}\label{sec:browser-talking-to-node}

Unlike an ISOTOPE node, an ACTOR node gives clients access to the full set of Erlang-style concurrency predicates and exposes them over a bidirectional WebSocket connection. This combination enables push-driven interaction, mixed-initiative communication, and concurrent actor programming -- capabilities that the HTTP-based APIs cannot support. Figure~\ref{fig:browser-to-toplevel} illustrates a scenario in which a user engages with an ACTOR node over such a connection.

%As indicated in the hierarchy of profiles in Figure~\ref{fig:node-profiles}, unlike an ISOTOPE node, an ACTOR node must give a client access to the full set of Erlang-style concurrency-oriented predicates listed in the diagram. Any type of actor can be employed, and a user is able to fully utilize its private database. 
%
%The bidirectional nature of the WebSocket API enables data to be \textit{pushed} from the node to the client. This functionality not only allows the ACTOR node to simulate input/output operations but also facilitates mixed-initiative interactions, where both the client and the node can actively participate in the communication.
%
%Most significantly, the WebSocket API unleashes the power of Erlang-style concurrent programming within the ACTOR node. By enabling actors to run concurrently and communicate seamlessly, the API supports highly dynamic and responsive interactions, empowering users to design complex, efficient, and flexible systems.
%
%Figure~\ref{fig:browser-to-toplevel} illustrates a scenario in which a user engages with an ACTOR node over a stateful WebSocket connection. 

\begin{figure}[h]
    \centering
	\includegraphics[width=\textwidth]{browser-to-actor-node}
    \caption{A scenario involving a web browser in conversation with a remote node running a toplevel actor over a stateful WebSocket connection.}
    \label{fig:browser-to-toplevel}
\end{figure}

After having loaded the shell and established a WebSocket connection with the node, the JavaScript process sent the following command to the node's controller, instructing it to spawn a toplevel actor running a session with the clause \texttt{q(X)\,:-\,p(X)} loaded into the actor's private database:

\begin{verbatim}
{ "command": "toplevel_spawn",
  "options": "[session(true), load_text('q(X):-p(X).')]" }
\end{verbatim}

\noindent As requested by the client, the pid of the toplevel actor was returned to the client, not in the form of a HTTP response, but embedded in a JSON formatted message sent over the WebSocket connection:

\begin{verbatim}
{ "type": "spawned",
  "pid" : 23981144 }
\end{verbatim}

\noindent From here on, the solving of the query \texttt{?-q(X)} proceeds in a way very similar to how it worked in the ISOTOPE scenario, except that instead of the JavaScript shell making an HTTP GET request, it sent the following command message to the node over the WebSocket connection:

\begin{verbatim}
{ "command": "toplevel_call",
  "pid": 23981144,
  "goal": "q(X)",
  "options": "[limit(1)]" }
\end{verbatim}

%\noindent The command told the node controller to make a call to \texttt{toplevel\_call/3} and to convert the resulting answer term into a JSON structure that encodes the first solution: 
%
%\begin{verbatim}
%{ "type": "success",
%  "pid" : 23981144,
%  "data": [{"X":"a"}],
%  "more": true }
%\end{verbatim}
%
%\noindent The value \texttt{true} of the \texttt{more} property signifies that more solutions may be available. Asking for the next solution results in a similar message, except that the value of the \texttt{data} property becomes \texttt{[\{"X":"b"\}]} and the value of \texttt{more} becomes \texttt{false}.

\noindent Knowing that the node's shared database offers a definition of \texttt{count\_actor/1}, the user then decided to spawn a count server of the kind we saw in Chapter~\ref{sec:web-prolog}. Here is the message sent over the WebSocket connection:

\begin{verbatim}
{ "command": "toplevel_call",
  "pid": 23981144,
  "goal": "spawn(count_actor(0),Pid),register(counter,Pid)",
  "options": "[limit(1)]" }
\end{verbatim}

\noindent This resulted in:

\begin{verbatim}
{ "type": "success",
  "pid" : 23981144,
  "data": [{"Pid":91883433}],
  "more": false }
\end{verbatim}

\noindent While the toplevel actor continued to serve shell queries, the spawned count actor ran concurrently and received \texttt{count} messages independently of the shell interaction, illustrating the concurrent, actor based nature of the ACTOR profile. With access to the pid of the spawned count server, the user could send it a \texttt{count} message in order to try it out, and call \texttt{flush/0} to inspect the response.

Finally, the user called \texttt{toplevel\_spawn/2} with the \texttt{node} option. This is yet another way to do network-transparent concurrency that will be dealt with in more detail in Chapter~\ref{sec:the-prolog-web}.

In this implementation, if the user leaves the terminal at \texttt{http://n7.org/shell}, closing the WebSocket causes the controller to terminate the session toplevel (via \texttt{exit/2}); before exiting, the toplevel terminates the spawned count actor. This yields a supervision-like hierarchy, discussed further in Chapter~\ref{sec:the-prolog-web}.


%Should the user leave the terminal at \texttt{http://n7.org/shell} the toplevel process will first kill (by calling \texttt{exit/2}) the count actor and then die (since the closing of the WebSocket will kill it, again by a call to \texttt{exit/2}). There is a kind of supervision hierarchy in play here, and we shall say more about that in Chapter~\ref{sec:the-prolog-web}.

\subsection{Social robot talking to an ACTOR node}\label{sec:robot-talking-to-actor-node}

In the previous scenarios, the shell and its toplevel actor served as the root of control: it was the first long-lived process from which other actors were spawned, monitored, and addressed. In general, however, the root of an application need not be a toplevel at all. Any actor can play that role, and a client may interact with such an actor directly, without routing interaction through a shell-style PTCP conversation.

Figure~\ref{fig:robot-to-actor-node} sketches a deliberately minimal toy scenario in which a social robot interacts with a remote ACTOR node over the stateful WebSocket API. 

\begin{figure}[h]
\centering
\includegraphics[width=11.5cm]{robot-to-actor-node}
\caption{Stateful interaction in a scenario involving a social robot and a remote ACTOR node.}
\label{fig:robot-to-actor-node}
\end{figure}

\noindent The robot itself provides speech recognition and speech synthesis: it turns the user's spoken input into text, sends that text to the node, receives a reply as text, and speaks it back to the user. In this toy setup the remote side does nothing clever -- it simply echoes the input -- so the interaction is a ``parrot'' loop. The purpose is not to propose a useful robot skill, but to isolate the transport-level point: unlike the shell scenarios, the client is not a terminal driving a toplevel, and yet it can still engage in persistent, bidirectional interaction with a remote node.

At the protocol level, the scenario begins with the robot requesting that the node spawn an echo actor. The request is sent as a WebSocket message to the node controller (encoded here in JSON for the benefit of a non-Prolog client), and the node responds by returning the pid of the spawned actor. This pid becomes the capability by which the robot can address the actor in subsequent messages.

\begin{verbatim}
{ "command": "spawn",
  "goal": "echo_actor",
  "options": "[monitor(true)]" }
\end{verbatim}

\begin{verbatim}
{ "type": "spawned",
  "pid": 47110023 }
\end{verbatim}

\noindent The WebSocket transport matters here because the interaction is push-driven and continuous. The robot establishes one connection and then exchanges messages in both directions as events occur: it sends new utterances as soon as they are recognised, and the node can return responses as soon as they are produced, without requiring a new HTTP request per response. This is the sense in which the interaction is stateful at the connection level: the channel remains open, latency is low, and messages can be delivered whenever either party chooses to speak.

In realistic conversational systems, of course, echoing recognised text is not enough. One typically needs additional components such as natural language understanding, a dialogue manager, grounding in a knowledge base, and response generation, as well as timing-sensitive coordination of speech, gaze, and facial expressions. We return to such agent-level concerns in Chapter~\ref{sec:prolog-ai-agents-on-the-agentic-web}; the present scenario only serves to show that, with the ACTOR profile, the transport and interaction model already supports the kind of persistent, low-latency, bidirectional exchange that embodied clients require.

%\subsection{Social robot talking to an ACTOR node}\label{sec:robot-talking-to-actor-node}
%
%In the previous scenarios, the shell and its toplevel actor served as the root of control: it was the first long-lived process from which other actors were spawned, monitored, and addressed. In general, however, the root of an application need not be a toplevel at all. Any actor can play that role, and a client may interact with such an actor directly, without routing interaction through a shell-style PTCP conversation.
%
%Figure~\ref{fig:robot-to-actor-node} sketches a scenario in which a social robot acts as a client of a remote ACTOR node over the stateful WebSocket API. The robot establishes a WebSocket connection to the node, spawns an application-specific actor, and then engages in a bidirectional exchange of Prolog terms that represent events (from the robot to the node) and actions (from the node to the robot). Unlike the shell scenarios, the spawned actor is not a toplevel actor but a \emph{dialog manager} — an actor that coordinates a conversation by reacting to incoming events, maintaining conversational state, and issuing actions such as speaking, asking follow-up questions, or performing non-verbal cues.
%
%\begin{figure}[h]
%\centering
%\includegraphics[width=11.5cm]{robot-to-actor-node}
%\caption{Stateful interaction in a scenario involving a social robot and a remote ACTOR node.}
%\label{fig:robot-to-actor-node}
%\end{figure}
%
%At the protocol level, the scenario begins with the robot requesting that the node spawn a dialog manager actor. The request is sent as a WebSocket message to the node controller (encoded here in JSON for the benefit of a non-Prolog client), and the node responds by returning the pid of the spawned actor. This pid becomes the capability by which the robot can address the actor in subsequent messages.
%
%\begin{verbatim}
%{ "command": "spawn",
%  "goal": "dialog_manager()",
%  "options": "[monitor(true)]" }
%\end{verbatim}
%
%\begin{verbatim}
%{ "type": "spawned",
%  "pid": 47110023 }
%\end{verbatim}
%
%\noindent Once the actor exists, the robot sends it events representing what the robot perceives or what the user has done. For example, a speech recogniser may produce a text hypothesis together with metadata such as a confidence value or timing information. The exact event vocabulary is application-specific, but the important point is that the interaction is phrased as the exchange of structured terms.
%
%\begin{verbatim}
%{ "command": "send",
%  "to": 47110023,
%  "msg": "user_utterance('hello', [confidence(0.92)])" }
%\end{verbatim}
%
%\noindent The dialog manager reacts by sending action terms back to the robot. Because the transport is stateful, the node can \emph{push} such actions as soon as they are produced, without waiting for the robot to poll for them. This is the key interactional difference compared to the HTTP-based APIs: output and other mid-computation responses become ordinary pushed messages rather than responses that must be collected by a long-polling-like pattern.
%
%\begin{verbatim}
%{ "type": "msg",
%  "from": 47110023,
%  "data": "say('Hello! How can I help?')" }
%\end{verbatim}
%
%\noindent In realistic conversational systems, timing is not an afterthought but a first-class constraint. The dialog manager may start timers (for turn-taking, barge-in, or silence detection), cancel pending actions when new events arrive, and interleave multiple streams of activity. These are precisely the kinds of behaviours that benefit from an actor-based design: timeouts can be realised by receive timeouts, delayed actions can be scheduled by spawning helper actors or by timer services, and independent concerns (speech, gesture, context updates) can be handled by separate actors that communicate asynchronously.
%
%\begin{verbatim}
%{ "type": "msg",
%  "from": 47110023,
%  "data": "set_timer(turn_timeout, 5000)" }
%\end{verbatim}
%
%\begin{verbatim}
%{ "type": "msg",
%  "from": 47110023,
%  "data": "gesture(smile)" }
%\end{verbatim}
%
%\noindent Conceptually, the dialog manager is well modelled as an event-driven state machine: it receives events, consults its private beliefs (for example in its private database), updates conversational state, and emits actions. This perspective aligns with the chapter's broader point about agents and protocols: what matters operationally is not that the client is a browser or a terminal, but that persistent agents exchange Prolog terms according to a protocol whose timing and ordering constraints are part of the behaviour. In the ACTOR profile, the WebSocket API exists to make this kind of low-latency, push-driven, interleaved interaction straightforward, while the actor model provides the concurrency substrate needed to keep such interaction responsive under competing demands.


\subsection{Could there be other profiles?}\label{sec:other-profiles}

The hierarchy RELATION $\subset$ ISOBASE $\subset$ ISOTOPE $\subset$ ACTOR is useful because many applications need less than a full ACTOR node. Simpler profiles are easier to implement, govern, and deploy, while still supporting portable client behaviour.

At the same time, the hierarchy is not meant to be a closed taxonomy. The existing profiles mostly specify \emph{what} capabilities are exposed at the node interface. One can also imagine profiles that express \emph{how} a node chooses to do it: restrictions, guarantees, or disciplines that cut across the capability levels.

A \textsc{PURE} profile is one example: clients would be restricted to a declarative subset (excluding cut, non-logical arithmetic, and related extra-logical constructs), while still potentially allowing features such as source injection. Another example is a statechart discipline exposed as an optional ACTOR sub-profile for reactive conversational control. Whether such orthogonal profiles are worth standardising remains open, but the architecture allows them.

%
%
%\subsection{Could there be other profiles?}\label{sec:other-profiles}
%
%The profile ladder is not meant to be a closed taxonomy. The existing profiles are primarily about \emph{what a node can do} at its interface. One can also imagine profiles that express \emph{how} a node chooses to do it: restrictions, guarantees, or disciplines that cut across the capability levels. A PURE profile, for instance, might admit only a completely pure subset of Prolog -- excluding cut, non-logical arithmetic, and other extra-logical constructs -- while still supporting source code injection. Such a profile would encourage a strictly declarative programming style and might simplify certain kinds of formal reasoning about distributed queries. Similarly, one could imagine a \textit{statechart discipline} exposed as a capability or sub-profile within ACTOR nodes: not a new rung in the hierarchy, but an optional, explicitly supported control layer for reactive, conversational behaviour. Whether the practical benefits of such orthogonal profiles would justify their additional specification and implementation complexity remains an open question, but the architecture does not foreclose the possibility.
%
%
%---
%
%The profile hierarchy reflects a key insight: many web applications need far less from a Prolog back-end than what a full ACTOR node provides. A simple ``logic server'' -- one that answers queries over stateless HTTP -- suffices for a wide range of applications. Since such a node is considerably easier to implement than an ACTOR node, the hierarchy lets implementors scale their commitment to what their use case actually requires.
%
%It is worth noting that the profile ladder is not meant to be a closed taxonomy. Looking ahead, one can also imagine additional profiles that are largely orthogonal to the current hierarchy. The existing profiles are primarily about \emph{what a node can do} at its interface. Orthogonal profiles would instead express \emph{how} a node chooses to do it: restrictions, guarantees, or disciplines that cut across the capability levels. A PURE profile, for instance, might admit only a completely pure subset of Prolog -- excluding cut, non-logical arithmetic, and other extra-logical constructs -- while still supporting source code injection. Such a profile would encourage a strictly declarative programming style and might simplify certain kinds of formal reasoning about distributed queries. Similarly, one could imagine a \textit{statechart discipline} exposed as a capability or sub-profile within ACTOR nodes: not a new rung in the hierarchy, but an optional, explicitly supported control layer for reactive, conversational behaviour. Whether the practical benefits of such orthogonal profiles would justify their additional specification and implementation complexity remains an open question, but the architecture does not foreclose the possibility.




%The profile hierarchy reflects a key insight: many web applications need far less from a Prolog back-end than what a full ACTOR node provides. A simple ``logic server'' -- one that answers queries over stateless HTTP -- suffices for a wide range of applications. Since such a node is considerably easier to implement than an ACTOR node, the hierarchy lets implementors scale their commitment to what their use case actually requires.
%
%Looking ahead, one could imagine additional profiles orthogonal to the current hierarchy. A PURE profile, for instance, might admit only the completely pure subset of Prolog -- excluding cut, non-pure arithmetic, and other extra-logical constructs -- while still supporting source code injection. Such a profile would encourage a strictly declarative programming style and might simplify certain kinds of formal reasoning about distributed queries. Whether the practical benefits would justify the additional complexity remains an open question, but the architecture does not foreclose the possibility.

\section{Notes on node implementation}\label{sec:node-implementation-notes}

This chapter specifies node behaviour primarily in terms of externally observable contracts: profiles describe what a client may rely upon, and the web APIs describe the shapes of requests and responses. Implementation techniques are not part of the contract, and different Prolog systems -- and different host platforms -- will naturally realise the contracts in different ways. Nevertheless, a few recurring engineering issues are worth highlighting. They explain several design choices, clarify why multiple APIs are worth supporting, and help readers anticipate the practical costs associated with each profile.

\subsection{Statelessness as a protocol property}\label{sec:statelessness-protocol}

The stateless HTTP API is \emph{stateless for the client}: each request contains all information needed to interpret it, and the node is not required to remember a conversational session in order to respond. This does not prohibit the node from keeping internal caches. It merely means that caches are best-effort optimisations whose presence must not be observable as a behavioural dependency. A node may evict cached state at any time and still remain correct and conformant.

As noted in Section~\ref{sec:browser-to-relation-node}, statelessness brings familiar web virtues: meaningful URIs, replayable requests, and cacheable responses. A corresponding limitation is that a stateless request cannot be \emph{externally} interrupted by the client once evaluation has begun. There is no control channel associated with the computation, and no session actor to signal. In practice, therefore, the stateless API must rely on node-enforced resource policy, most importantly running timeouts, to avoid nonterminating or overly expensive computations.

\subsubsection{Two strategies for implementing paged solution retrieval}\label{sec:two-strategies-paged-solution-retrieval}

The \texttt{offset} and \texttt{limit} parameters make it possible to retrieve solutions in slices. This subsection is concerned only with how a node might implement that contract efficiently; the architectural interpretation of paged retrieval as an instance of Prolog-style backtracking is discussed separately in Chapter~\ref{sec:the-prolog-web}.

A naive implementation uses \emph{restart-and-skip}: to compute the slice starting at offset $N$, the node restarts evaluation from the beginning, discards the first $N$ solutions, and then collects up to \texttt{limit} further solutions. This strategy is simple and portable, but it may waste work when later slices are requested, because earlier solutions are recomputed only to be thrown away.

A more efficient strategy uses \emph{save-and-resume}. The node evaluates the goal in an execution context that can be suspended after producing a slice and resumed later to continue search. The node may maintain a cache keyed by a fingerprint of the query (goal, template, and any relevant execution parameters) together with the next slice boundary. When a request arrives for offset $N$, the node first attempts to resume a cached computation that has already produced $N$ solutions; if none is available, it falls back to restart-and-skip. This improves performance in the common case without changing the protocol contract.

Importantly, the choice between these strategies is not visible at the API level. Clients only observe that slices are returned in Prolog order and that the \texttt{more} flag correctly reports whether further solutions may exist. All optimisation is therefore constrained by the requirement that cached state is disposable: eviction must never change the meaning of subsequent responses, only their cost.


\subsection{Semi-stateful HTTP sessions and response retrieval}\label{sec:semi-stateful-notes}

The semi-stateful HTTP API differs from the stateless API in one crucial respect: each request targets a long-lived toplevel actor, identified by a pid, and the observable behaviour of the conversation depends on the evolving state of that actor. Nevertheless, each individual HTTP exchange is still an ordinary request--response transaction. In this sense, the protocol is \emph{semi-stateful}: the transport remains stateless, but the addressed computation is not.

From an implementation perspective, this style aligns naturally with the architecture of SWI-Prolog's \texttt{library(pengines)} and systems such as SWISH: the server maintains per-session engines, while the client repeatedly issues short requests that either advance the computation, or retrieve the next available response produced by that computation. 

Output and prompts fit this model once they are treated as ordinary protocol events emitted by the session actor, rather than as exceptional side effects. Concretely, the toplevel actor can be understood as producing a sequence of \emph{responses} (answers, failure, errors, output fragments, prompts, and acknowledgements) that the client retrieves over a series of short HTTP exchanges.

\begin{figure}[h]
    \centering
	\includegraphics[width=10.5cm]{PTCP-HTTP.pdf}
    \vspace{0.4cm}\caption{Statechart specifying the PTCP for a successful conversation with a toplevel over HTTP. The transitions are labeled with \textit{message types}. Types in bold are sent from the client to the toplevel, whereas message types with a leading \texttt{/} goes in the opposite direction, from the toplevel to the client.}
    \label{fig:statechart-3}
\end{figure}


The statechart in Figure~\ref{fig:statechart-3} extends the core PTCP with the states and transitions needed for HTTP delivery. When the toplevel emits an \texttt{/output} event during evaluation, the client enters a local pull cycle: it retrieves the pending event and immediately issues a follow-up \textbf{poll} request for the next pending response. This cycle repeats until the toplevel produces a terminal response (\texttt{/success}, \texttt{/failure}, \texttt{/error}), at which point control returns to the main protocol states. The mechanism resembles long-polling only locally -- within the scope of a single goal evaluation that happens to produce intermediate events -- rather than as a general-purpose polling loop.

Prompts require one additional step: when the next pending response is a prompt, the client retrieves it, presents it to the user, and sends a follow-up request that delivers the user's reply back to the session actor. From the protocol's perspective, this is still a request--response rhythm and no polling is needed here; the client is simply advancing the conversation by alternately retrieving pending events and, when prompted, supplying input. Ordering is preserved because responses are dequeued in the order they were produced. Transient transport failures may delay retrieval but do not change the logical meaning of the session.
%Ordering is preserved because responses are dequeued in the order they were produced, and robustness follows from the same principle as before: transient transport failures may delay retrieval, but must not change the logical meaning of the session -- only its latency.

%A practical way to realise this over plain HTTP is to keep at most one outstanding ``await response'' request per session pid. After the client sends a PTCP step (for example \texttt{call}, \texttt{next}, or \texttt{respond}), it asks for the next response event. If an event is already queued, the node returns it immediately; otherwise, the node holds that request briefly and returns as soon as one event is produced (or a timeout occurs). If more events are expected (for example another \texttt{output}), the client issues a new await request. In this sense, the long-polling is local and demand-driven, while still allowing prompts and streamed output during evaluation.


%A practical way to realise this over plain HTTP is a pull-based event stream. The client repeatedly issues a ``next response'' request against the session pid. If no response is currently available, the server holds the request for up to some timeout (a long-polling pattern), and returns as soon as the actor emits the next event. The client then immediately issues a new request to await the next response. This approach works well with browser constraints and common HTTP infrastructure, while still allowing the session actor to publish output and prompts during evaluation.



%\todo{The material below needs more work.}
%
%Output and prompts fit this model by being treated as responses that may occur mid-computation rather than as exceptional side effects.
%
%In practice, response retrieval over HTTP is often realised using a long-polling-like pattern: the client issues a request that blocks until the next response becomes available (success, failure, error, output, prompt, abort acknowledgement), then immediately issues a new request to await the next one. This preserves compatibility with HTTP infrastructure and browser constraints while still allowing the session actor to emit output and prompts during evaluation.

\subsection{Stateful WebSockets and push-driven interaction}\label{sec:websocket-notes}

The stateful WebSocket API exists for interactions that are naturally push-driven. Once a connection is established, the node can deliver answers, output fragments, prompts, and monitor events as soon as they are produced, without requiring a new client request for each response. This improves latency and avoids the request choreography of our long-polling trick. It also supports interleaved traffic from multiple actors over one connection, which is central to ACTOR-profile use cases.

Implementation complexity, however, shifts from request handling to connection lifecycle management. A robust node must maintain per-connection routing state (which pids belong to which client), preserve ordering within each actor stream, and enforce outbound back-pressure so that a slow client cannot cause unbounded buffering on the node. It must also define clear disconnect semantics: when a client vanishes, session actors and child actors are either terminated or detached according to policy, and monitored parties observe the corresponding \texttt{down}/\texttt{exit} events. In this sense, the WebSocket API is conceptually simple but operationally demanding: the protocol surface is small, while lifecycle and resource handling must be precise.



%\subsection{Stateful WebSockets and push-driven interaction}\label{sec:websocket-notes}
%
%The stateful WebSocket API primarily exists to support bidirectional, low-latency interaction and true push semantics. Unlike HTTP, a WebSocket connection allows the node to deliver responses as soon as they are produced, without requiring the client to issue a new request. This matters both for responsiveness (for example, streaming output) and for actor-style programming, where multiple actors may send messages to a client in an interleaved fashion.
%
%Implementation-wise, this shifts complexity away from request scheduling and towards connection management: framing and parsing messages, maintaining per-connection routing state, enforcing per-connection resource limits, handling disconnects, and ensuring that actor termination and monitoring semantics remain coherent when the client disappears. In other words, the WebSocket API is not conceptually deep, but it is operationally demanding: it requires robust server and client libraries, and careful integration between the node controller, actor lifecycles, and connection lifecycles.

%\subsection{Resource bounds, eviction, and safety}\label{sec:implementation-resource-safety}
%
%Save-and-resume for the stateless API shifts cost from CPU time to memory and resource management. Semi-stateful sessions likewise require that the node holds conversational state for as long as a pid remains valid, and the WebSocket API adds the operational cost of long-lived connections. A practical node therefore requires resource bounds, eviction policies, and admission control to prevent exhaustion.
%
%These considerations interact with profile obligations. Less capable profiles are easier to implement safely because they admit fewer long-lived resources. More capable profiles -- notably those that admit long-lived sessions, dynamic database updates, and arbitrary client-supplied code -- require correspondingly stronger governance mechanisms: running and idle timeouts, caps on injected code size and private database size, limits on concurrent actors and sessions per client, and policies for reclaiming abandoned computations. %Appendix~\ref{sec:profiles-and-implementation} sketches one concrete realisation of these ideas.
%
%A recurring design theme is that the protocols are specified so that resource management remains a matter of policy rather than semantics. Caches may be dropped, sessions may expire, and connections may be closed under clear rules, without forcing changes to the meaning of the programming model. The goal is to keep interoperability grounded in observable behaviour, while leaving implementers free to choose the engineering techniques that best fit their platform.

%\section{Notes on node implementation}\label{sec:node-implementation-notes}
%
%This chapter specifies node behaviour primarily in terms of externally observable contracts: profiles describe what a client may rely upon, and the web APIs describe the shapes of requests and responses. Implementation techniques are not part of the contract, and different Prolog systems -- and different host platforms -- will naturally realise the contracts in different ways. Nevertheless, a few recurring engineering issues are worth highlighting, because they explain several design choices and motivate the existence of multiple APIs.
%
%\subsection{Statelessness as a protocol property}\label{sec:statelessness-protocol}
%The stateless HTTP API is \emph{stateless for the client}: each request contains all information needed to interpret it, and the server is not required to remember a conversational session in order to respond. This does not prohibit the node from keeping internal caches. It merely means that caches are best-effort optimisations whose presence must not be observable as a behavioural dependency. A conformant node may evict cached state at any time and still remain correct.
%
%\subsection{Two strategies for implementing paged solution retrieval}\label{sec:paging-two-strategies}
%The \texttt{offset} and \texttt{limit} parameters make it possible to retrieve solutions in slices. A naive implementation uses \emph{restart-and-skip}: to compute the slice starting at offset $N$, the node restarts evaluation from the beginning, discards the first $N$ solutions, and then collects up to \texttt{limit} further solutions. This strategy is simple and portable, but it may waste work when later slices are requested, because earlier solutions are recomputed only to be thrown away.
%
%A more efficient strategy uses \emph{save-and-resume}. The node evaluates the goal in an execution context that can be suspended after producing a slice and resumed later to continue search. The node may maintain a cache keyed by a fingerprint of the query (goal, template, and any relevant execution parameters) together with the next slice boundary. When a request arrives for offset $N$, the node first attempts to resume a cached computation that has already produced $N$ solutions; if none is available, it falls back to restart-and-skip. This improves performance in the common case without changing the protocol contract.
%
%\subsection{Resource bounds, eviction, and safety}\label{sec:implementation-resource-safety}
%Save-and-resume shifts cost from CPU time to memory and resource management. A practical node therefore requires eviction policies, timeouts for suspended computations, and admission control to prevent resource exhaustion. These considerations interact with profile obligations: less capable profiles are easier to implement safely, while more capable profiles (notably those that admit long-lived sessions and arbitrary user code) require correspondingly stronger governance mechanisms. Appendix~\ref{sec:profiles-and-implementation} sketches one concrete realisation of these ideas.

\section{Summary}

Chapter~\ref{sec:web-prolog} equipped Web Prolog with actors and concurrency, and this chapter lifted those ideas to the level of Prolog agents and nodes. The central move was to treat toplevels as stateful actors governed by an explicit Prolog Toplevel Communication Protocol (PTCP). By making the conversational contract first-class -- including how queries are started, how solutions are streamed, how I/O and aborts are handled, and how sessions are terminated -- we turned an informal REPL interaction pattern into a well-defined behavior that can be implemented inside a node or exposed over the network.

On top of this, we introduced Prolog nodes as agents in their own right. A node is more than a mere container for code and data: it offers one or more toplevels, may host arbitrary actors, and is addressable on the Prolog Web through URIs and web APIs. The notion of node profiles (RELATION, ISOBASE, ISOTOPE, ACTOR) makes explicit which capabilities a node must offer to its clients. In particular, the ACTOR profile marks nodes that support the full actor model and PTCP-based toplevels, whereas the more restricted profiles support progressively smaller subsets of those capabilities. This lets implementors scale their commitment, and lets clients reason about what they can rely upon in a heterogeneous network.

A recurring theme in the chapter was the distinction between private state and shared data. Actors and toplevels maintain their own conversational and internal state -- their ``private beliefs'' -- while nodes may expose shared knowledge bases that serve as common background for all actors running on that node. The dynamic database, though available within an actor, remains process-local and cannot be used for inter-process communication, preserving the principle that processes should communicate to share memory rather than share memory to communicate.

The three communication paradigms -- stateless HTTP, semi-stateful HTTP, and stateful WebSocket -- serve complementary purposes. Stateless interactions offer simplicity, cacheability, and horizontal scalability; semi-stateful sessions enable dynamic database updates without persistent connections; stateful WebSockets unlock the full power of bidirectional, real-time interaction and concurrent actor programming. Supporting all three within a single node architecture allows applications to choose the appropriate trade-off between simplicity and capability.

Finally, the chapter touched on security considerations. Like JavaScript, Web Prolog is designed as a sandboxed language: it omits file I/O, socket programming, and direct OS access, relying instead on the host environment for such capabilities. Security in the actor model rests primarily on name distribution -- if a process identifier is never revealed, the process cannot be accessed. This principle, combined with the constrained language subset and profile-based capability advertisement, forms the foundation of a security model appropriate for running untrusted code on shared infrastructure.

The overall message of the chapter is that Web Prolog is not just a language but also a way of packaging processes, protocols, and capabilities into networked agents that can be composed and reasoned about at a higher level of abstraction. Chapter~\ref{sec:the-prolog-web} extends these local concepts across the network, showing how nodes discover each other, how actors migrate and communicate across node boundaries, and how the Prolog Web emerges as a federated infrastructure for distributed logic programming.

\medskip
\noindent The motivation for profiles is not only security, but also coordination. FYPB emphasises that portability is not just about language constructs, but also about portable libraries and developer tools -- for unit testing, documentation, linting, and package distribution. A profile is a way to make that expectation concrete: it specifies which predicates, libraries, and APIs are guaranteed to exist on a node of a given kind. In other words, profiles separate \emph{observable capability} from \emph{implementation commitment}, enabling tools to target stable capability sets even when different implementations continue to innovate internally.

%\section{SCRAP}
%
%
%\subsection{Actors, toplevels, and the PTCP}
%We unify the treatment of actors and toplevel actors by making explicit how the Prolog Toplevel Communication Protocol (PTCP) governs interactions. Toplevels are stateful actors whose conversational state is part of the protocol contract; this frames what nodes must expose to clients and what clients may rely upon. Briefly restate the PTCP guarantees and the role of the toplevel as a first-class actor, then summarize how these micro-semantics underpin the remainder of the chapter.
%\emph{Implication.} Keep PTCP semantics close to the implementation surface and reference them whenever interface obligations or timeouts are discussed.
%
%---
%
%The role of toplevel actors is introduced in Section~\ref{sec:toplevel}, and for the protocol and state-machine details, see Section~\ref{sec:pcp-protocol}.
%
%
%
%
%\subsection{Nodes as agents: capabilities, profiles, and conformance}\label{sec:conformance}
%Treat a node as an agent with advertised capabilities. Connect those capabilities to \emph{profiles} that specify language features and interaction affordances. Explain conformance as the checkable link between what a node claims to support and what clients may assume. Retain the short forward-look on additional profiles (including a possible PURE profile) as a closing paragraph of this subsection.
%\emph{Implication.} Profiles are the stable contract for portability; nodes may under-advertise but must not over-claim.
%
%---
%
%The notion of a node as an agent is developed in Section~\ref{sec:nodes-as-agents}, and the profile system is defined in Section~\ref{sec:node-profiles}. Conformance and advertised capabilities are revisited in Section~\ref{sec:conformance}.
%
%
%
%
%
%\subsubsection{Some scenarios}
%
%\noindent These options are particularly useful in the following scenarios:
%
%\begin{itemize}
%
%\item Pagination of results: When dealing with a large dataset, it is often impractical to retrieve and display all the records at once, especially in web applications. Instead, results are divided into pages, with each page showing a smaller subset of the data.
%
%\item \textit{Next} and \textit{Previous} controls.
%
%\item Efficient data retrieval: In cases where only a portion of the dataset is needed for processing, using \texttt{offset} and \texttt{limit} helps in efficiently retrieving just that portion, reducing the load on the node and the network.
%
%\item Implementing infinite scrolling: In user interfaces that implement infinite scrolling (such as social media feeds), \texttt{offset} and \texttt{limit} can be used to load more data dynamically as the user scrolls down the page.
%
%\item Handle an infinite number of solutions: 
%
%\item Limit the number of network roundtrips on the Prolog Web
%
%\item Avoiding recomputation:
%
%\end{itemize}
%
%\noindent A special case of pagination is when we choose to set \texttt{limit} to \texttt{1} in order to make it much easier to implement a traditional Prolog terminal in JavaScript.
%
%\subsubsection{Do we need more profiles?}\label{sec:more-profiles}
%
%A possible advantage with the ISOBASE and ISOTOPE profiles is that they, due to the lack of operations for manipulating the dynamic database (in ISOBASE) and for general I/O (in both), encourage a programming style focusing on the declarative, logical aspects of Prolog. Those profiles do not rule out impure constructs altogether though, since the cut and non-pure arithmetic are still available.
%
%To encourage an even greater focus on declarative programming techniques, we could imagine a profile -- perhaps to be named PURE? -- that allows only a completely pure subset of Prolog to be used in queries and programs, thus leaving out constructs such as the cut and non-pure arithmetic. We note that the capabilities of such a profile would be orthogonal to the capabilities of the ISOBASE profile. In particular, we see no reason why a pure profile should not be able to support source code injection. 
%
%\subsubsection{Motivating the profile hierarchy}\label{sec:profiles-motivations}
%
%A key insight here is that many web applications can be written with much less support from the Prolog back-end than what is offered by a node implementing the ACTOR profile. Indeed, many web applications are written in HTML, CSS and JavaScript and have only need for what amounts to a ``logic server". Since a node only capable of acting as a logic server is a lot easier to implement than an ACTOR node, this should send us off in the right direction.
%
%On the other hand, there are applications that need the extras provided by the ACTOR profile.
%
%
%---
%
%Another (but less important) reason for defining an hierarchy of profiles is that it allows the owner of a node to \textit{disguise} it as a node with lesser capabilities, and thus to offer clients \textit{less} than the node is actually capable of. As long as ``important'' clients are happy with what is on offer, the node owner may lack incentive for enabling other clients in the wild to run Web Prolog processes on the node. The idea is to enable the owner of (say) the most capable node to change a setting so that the node will present itself as less capable than it actually is. The motive for doing so may vary, but taking control over the management of resources (i.e. CPU cycles and memory capacity) might be one of them, security concerns another.
%
%At the same time, the \textit{owner} of the node has access to the capabilities if the full ACTOR profile. \todo{More here!}
%
%\subsubsection{Why so many communication models?}
%
%Supporting three communication paradigms -- stateless over HTTP, semi-stateful over HTTP, and stateful over WebSockets -- offers a node the flexibility to meet diverse client and application requirements while optimizing for performance, scalability, and usability. Each paradigm serves a unique purpose, addressing specific interaction models and constraints. Together, they create a robust architecture capable of handling the demands of modern web applications.
%
%Stateless communication over HTTP represents the simplest and most scalable paradigm. It aligns with the REST architectural style, where each request is independent, containing all the necessary information for the node to process it. This design minimizes node-side complexity, enabling horizontal scaling without session management overhead. Stateless APIs are ideal for high-traffic applications, where simplicity and efficiency are paramount, such as public-facing APIs, data retrieval services, or systems where caching plays a significant role. However, stateless interactions are limited in their ability to provide contextualized responses, as the node lacks memory of prior interactions.
%
%Semi-stateful communication over HTTP addresses the limitations of statelessness by introducing mechanisms to retain a lightweight session context. Through using a pid as a session identifier, the node can maintain a degree of continuity across requests. The semi-stateful API strikes a balance between simplicity and functionality, and is therefore well-suited for applications that require moderate personalization without the complexity of a fully stateful connection, such as e-commerce platforms, dashboards, or applications with short-lived user interactions.
%
%Stateful communication over WebSockets represents the most advanced but also the most resource-intensive paradigm. Unlike HTTP's request-response model, WebSockets establish a persistent, bidirectional connection between the client and node. This allows for real-time data exchange with minimal latency, making WebSockets indispensable for applications requiring immediacy, such as collaborative tools, conversational systems, or multiplayer games. The stateful nature of WebSockets enables a continuous dialogue between client and node, but it also demands careful resource management. Each connection consumes memory and computational resources, requiring the node to balance scalability with responsiveness.
%
%Integrating these three paradigms within a single node architecture allows for the optimal handling of diverse workloads. Stateless APIs can handle high-volume, low-complexity tasks, semi-stateful APIs can manage user-centric workflows, and stateful APIs can drive real-time interactivity. This modular approach not only enhances the node's versatility but also ensures that each communication style is used appropriately based on its strengths and limitations.
%
%Moreover, supporting all three paradigms fosters a seamless developer experience. Developers can choose the most suitable API style for their use case without being constrained by the node's capabilities. A unified authentication framework, shared data models, and consistent endpoint design further simplify the integration of these paradigms into client applications.
%
%The ability to support multiple communication paradigms also future-proofs the node, accommodating the evolving needs of applications and users. As the web ecosystem continues to expand, the demand for real-time interactions, personalized services, and scalable solutions will only grow. A node that can seamlessly transition between stateless, semi-stateful, and stateful interactions positions itself as a versatile backbone for modern and emerging technologies. %\todo{Add something about the transport option and the possibility to add other transports such as tcp. Also, perhaps move to chapter 4? }
%
%---
%
%Bi-directionality and statefulness are properties that make WebSocket by far the most flexible transport protocol for the Prolog Web. Over a WebSocket connection, a client can be in almost total control of an actor such as a toplevel. (We write ``almost'' here because the owner of the node on which the actor is running always has the ultimate say when it comes to how much resources will be allocated to the running of an actor.) 
%
%\todo{Say something about SWISH and library(pengines) and ite implementation of the semi-stateful HTTP API. Problem is that it tries to do too much.}
%
%\subsection{State and knowledge: private beliefs vs shared data}
%Consolidate the story of state in one place. Distinguish an actor's private dynamic database (its beliefs) from shared node-resident data (common knowledge). Discuss trade-offs: isolation, reproducibility, and performance against sharing, caching, and coordination. Clarify where persistence hooks may attach without violating actor encapsulation.
%\emph{Implication.} Treat private state as volatile and local; treat shared data as governed by explicit node policies.
%
%---
%
%Private per-actor dynamic databases are defined in  Section~\ref{sec:dynamic-private-database}, while node-resident shared data is introduced in Section~\ref{sec:shared-db}. For where persistence hooks and policies attach without breaking actor encapsulation, compare Section~\ref{sec:nodes-as-agents} and Section~\ref{sec:node-resident-actors}.
%
%---
%
%There are private static databases and private dynamic databases. There is a static shared database, but there is \textit{no} dynamic shared database, at least not for clients.
%
%
%
%
%
%%---
%%
%%\begin{enumerate}
%%
%%\item A \textit{Prolog agent} is a software process capable of communication and other forms of continuous interaction with other Prolog agents on the Prolog Web. As their language of communication they use Prolog terms of arbitrary complexity and even (parts of) whole Web Prolog programs.
%%
%%\item Prolog agents come in different types -- Prolog actors, Prolog toplevels and statechart actors. Some are short-lived, others are long-lived. Some are simple, others complex.
%%
%%\item Web Prolog is a web agent programming language, allowing us to program the agents as well as query them.
%%
%%\item We have demonstrated that Prolog-style backtracking search and Erlang-style concurrency can be integrated in a clear and elegant way if the receive predicate is allowed to be semi-deterministic, and we have argued that this is a natural way to supply nondeterministic querying with a deterministic interface that works well also in a multi-actor environment. %, and to fail only if the body of a receive clause fails. 
%%
%%\item This allowed us to encapsulate the behavior of a Prolog toplevel as a 
%%programming abstraction built on top of an actor running a suitable protocol. toplevel actors are first class citizens in the sense that new ones can be spawned and they (i.e. their pids) can be passed around in arguments. A Prolog shell is a Prolog toplevel equipped with facilities for interactive programming from a terminal.
%%
%%\item We have shown that the actor programming model of Web Prolog is able to co-exist harmoniously with its logic programming model, and that both models are usable in the same program, with actors used for modelling the reactive, deterministic interfaces to the user and other external entities, and with the nondeterministic component behaving as an actor, encapsulating the nondeterminism. 
%%
%%\item Also, it makes sense to think of a node as a kind of agent. 
%%
%%\item A node has a dual identity, as a Web Prolog runtime system, and as a node in the network forming the Prolog Web, in many ways similar to an ordinary web server. We have shown, that nodes come with different profiles with varying capabilities. 
%%
%%\end{enumerate}
%
%\subsection{Extending the capabilities of a node}
%
%A node may have capabilities well beyond what is offered by Web Prolog. 
%
%---
%
%In this way, the node can offer built-in predicates that provide functionalities not provided Web Prolog, such as support for constraint logic programming, allowing the use of constraints (e.g., arithmetic, finite domains) in addition to standard logical inference. However, extensions in the form of new data types, such as dictionaries, arrays, or graphs is not possible.
%
%---
%
%SICStus manual: For an introduction and a deeper description, see [Carreiro \& Gelernter 89a] or [Carreiro \& Gelernter 89b], respectively.
%
%\begin{verbatim}
%:- use_module(linda, [out/1, rd/1, in/1]).
%\end{verbatim}
%
%%\subsection{Implementing a node}\label{sec:implementing-node}
%%
%%Appendix~\ref{sec:profiles-and-implementation} describes how to implement a Prolog node by wrapping networking code around an existing Prolog system. The process involves implementing a stateless HTTP API, a version of nondeterministic remote procedure calls on top of the API, and a runnable specification for how predicates such as \texttt{spawn/1-3}, \texttt{!/2}, and \texttt{receive/1-2} should work.
%%
%%The implementation of the Prolog toplevel behavior and the built-in predicates for controlling it is also discussed. The implementations are partial, but the author indicates what else would be needed to complete them.
%%
%%---
%%
%%The web server that implements SWISH on top of \texttt{library(pengine)} takes us fairly close to an ISOTOPE node.
%%
%%---
%%
%%In the interest of making it as easy as possible to implement an ISOBASE node, we propose that \texttt{op/3}, the ISO Prolog predicate allowing a programmer to define prefix, postfix or infix operators and specify their precedences, does not have to be supported. After all, it is normally used as a directive, so to offer it in a profile that does not also support the loading of data and programs into the private database of an actor does not make much sense. Presumably, not having to support \texttt{op/3} should make it fairly straightforward to implement a parser for Web Prolog, or at least easier than if \texttt{op/3} would have to be supported.
%
%
%\subsection{Application patterns: from scripting to services}
%
%Show how the pieces come together for developers: (i) lightweight browser-side scripting against toplevels; (ii) node-resident services that use actors, supervision, and shared predicates; (iii) guidance on when to choose which profile and API. Use a compact example or two to illustrate the progression from script to service.
%\emph{Implication.} The same conceptual core scales from ad hoc scripts to durable services by tightening profiles, APIs, and supervision.
%
%---
%
%Browser-side scripting against toplevels is exemplified in Section~\ref{sec:web-prolog-as-a-scripting-language}; node-resident services and supervision are introduced in Section~\ref{sec:node-resident-actors}. Guidance on choosing a profile for each pattern aligns with Section~\ref{sec:node-profiles}.
%
%---
%
%\subsubsection{Building simple web applications in Web Prolog}\label{sec:simple-web-apps}
%
%A \emph{playground} for a programming language typically refers to an
%interactive, user-friendly environment that allows developers to write and execute code snippets in that language. Playgrounds are particularly valuable for educational purposes, exploratory programming, or testing smaller features in isolation without the need for full development environments.
%
%Figure~\ref{fig:browser-node-interaction-2} shows an application in a browser window that captures we might think of as the \textit{essence} of a playground for Web Prolog.
%
%\begin{figure*}[h]
%    \centering
%	\includegraphics[width=11cm]{wp-playground.png}
%    \caption{The essence of a Web Prolog playground.}
%    \label{fig:browser-node-interaction-2}
%\end{figure*}
%
%\noindent A playground for full Web Prolog must either be run in browser-based runtime environment, or talk to a node over the WebSocket sub-protocol. Note that the stateless HTTP API is not sufficient.
%
%Our GUI also features an editor, it means that when it is dirty, we can terminate the current toplevel and spawn a new one.
%
%\begin{verbatim}  
%connection.send(JSON.stringify({
%    command:"toplevel_spawn",
%    options:"[load_text('" + editor.content + "')]"
%}))
%\end{verbatim}  
%
%\noindent 
%
%It deserves to be pointed out that a playground like this obviously also could offer playing Erlang-style concurrent and distributed programming.  
%
%--- 
%
%Compute server - typically no owner-defined predicates, no shared database.
%
%---
%
%Real-Time Applications: WebSockets excel in real-time communication scenarios, such as chat applications, live gaming, or collaborative platforms, where continuous data needs to be exchanged between the client and server.
%
%Low-Latency Data Streams: WebSockets are ideal for sending real-time, low-latency data streams (e.g., stock prices, live sports scores, or multiplayer game state updates).
%
%---
%
%Figure~\ref{fig:browser-node-interaction-3} shows two applications that might be powered by Web Prolog.
%
%\begin{figure*}[h]
%    \centering
%	\includegraphics[width=11.5cm]{wn-tabling.png}
%    \caption{Two other applications built using Web Prolog technology}
%    \label{fig:browser-node-interaction-3}
%\end{figure*}
%
%\noindent The app in the back allows a user to explore the relation between grammars, strings and parse trees by means of a tiny interpreter for a subset of Definite Clause Grammar (DCG) running on the server and JavaScript tree drawing routines running on the client.
%
%The app in the front allowing a user to search the WordNet lexical dabase for word forms using a *-pattern and an optional filter on part of speech. 
%
%---
%
%Wordnet tabling not suitable for Prolog in browser.
%
%They both uses stateless HTTP API since stateful API not necessary.
%
%Apps easy to secure.
%
%---
%
%We must admit that looking back at the elevator pitch at the start of this chapter, it is clear that it made promises around Prolog agents that this chapter has not fulfilled. However, more exciting applications will be described in Chapter~\ref{sec:building-applications}.
%
%---
%
%\todo{Comment on remarks on UI development in FYPB}
%
%\subsubsection{Actors and actor organizations}\label{sec:actor-org}
%
%To is not at all unusual to consider an organization such as a company a kind of agent. An organization can act, it can talk to you. 
%
%They have the shared knowledge in common, etc.
%


%\section{SCRAP: Node profiles -- version with requirements}\label{sec:node-profiles}
%
%Not all Prolog nodes are equal; they vary in the capabilities they offer clients. We call the set of capabilities that a node provides its \emph{profile}. In this chapter we distinguish between two broad kinds of capabilities: \textit{language capabilities} (what Prolog fragments and predicates are available) and \textit{interactional capabilities} (what transports and web APIs the node supports). Four profiles will be presented and discussed in some detail: RELATION, ISOBASE, ISOTOPE and ACTOR. The Venn diagram in Figure~\ref{fig:node-profiles} provides an overview of how these profiles are related in our proposed design of the Prolog Trinity ecosystem. 
%
%\begin{figure}[h]
%    \centering
%	\includegraphics[width=10cm]{profile-hierarchy-new-2}
%    \vspace{0.3cm}\caption{A hierarchy of node profiles.}
%    \label{fig:node-profiles}
%\end{figure}
%
%\noindent Intuitively, ACTOR $\supset$ ISOTOPE $\supset$ ISOBASE $\supset$ RELATION: every ACTOR node is also an ISOTOPE node, and so on, but not vice versa. The RELATION profile is the least capable of the four.
%
%We borrow the term \emph{profile} in the standards sense: a tailored subset that constrains and organises a fuller design to meet particular goals (simplicity, performance, security, portability). Less capable profiles are easier to specify, implement, and standardise -- while still participating in the same ecosystem.  
%
%For example, owner-defined shared predicates and the HTTP API are located inside the area marked RELATION to signify that these features are supported by \emph{all} profiles, whereas placing the WebSocket API in the outermost area indicates that it is only available in the ACTOR profile.
%
%In short, when running against a node that offers the ACTOR or ISOTOPE profile, we can both query and program it. We may load predicates into a toplevel at the time of its creation, or use \texttt{assert/1} and \texttt{retract/1} to update its private database at any time during its lifecycle. Only the ACTOR profile, however, allows us to use predicates such as \texttt{spawn/3}, \texttt{send/2} and \texttt{receive/1} in queries and programs. Running against an ISOBASE or RELATION node, we can only query it, and of these two only the ISOBASE profile allows us to pose queries using built-in predicates. 
%
%We distinguish between three levels. First, a node implementation may support a rich set of capabilities in principle. Second, a particular deployment of that implementation may be configured to enable only a subset of these capabilities. Third, a running node instance advertises a profile to clients: the capabilities that clients may rely on. A node implementation may support more than the advertised profile, but it obviously cannot offer clients capabilities it does not implement.
%
%To be considered \textit{conformant}, a node \textit{must} provide authorized clients with \textit{all} capabilities specified by its profile. For instance, an ACTOR node \textit{must} offer clients access to the three web APIs and the entire set of built-in predicates indicated in Figure~\ref{fig:node-profiles}. Similarly, an ISOTOPE node \textit{must} provide clients with the two HTTP APIs and all predicates defined in the ISOTOPE subset of Web Prolog. Other profiles follow analogous requirements.
%
%A node \textit{may} provide \textit{more} than what is mandated by its profile, including additional predicates or web APIs. However, clients must not depend on the availability of these extra features. For example, a RELATION node may support the \texttt{rpc/2-3} predicate as well as both the stateless and semi-stateful HTTP APIs, but it is not obligated to do so. Only the stateless API is required.
%
%The following subsections present the requirements that each node profile must adhere to, accompanied by scenarios illustrating interactions between end-user clients, nodes, and actors. These scenarios aim to demonstrate the various types of communication that can occur, starting with the loading of a Prolog shell -- a specialized terminal implemented in JavaScript and executed within a web browser.
%
%\subsection{Browser talking to a RELATION node}\label{}
%
%\subsubsection{Profile requirements}
%A conformant RELATION node \textit{must} offer clients a stateless HTTP API for querying owner-defined shared predicates. At minimum, it must provide an endpoint for issuing queries (such as \texttt{/call}) that supports goal arguments and answer pagination (for example via \texttt{limit} and \texttt{offset} parameters) and returns results in a well defined, machine readable format such as JSON or Prolog terms. Viewed from the outside, a RELATION node must behave as if it were a pure relational Prolog server: it must support the discovery of zero, one or many solutions and allow clients to iterate lazily over these solutions, even if it implements this behaviour without genuine backtracking internally.
%
%A RELATION node is not required to make any ISO Prolog built-in predicates or control-constructs available to clients, not even constructs such as \texttt{,/2} or \texttt{;/2}. It must nevertheless distinguish between successful queries, failing queries and erroneous queries, and encode these outcomes explicitly in its responses.
%
%A RELATION node \textit{may} provide more capabilities than this minimum, such as additional HTTP endpoints, richer response formats or access to selected built-in predicates, but clients must not depend on such extensions. Any node that fails to provide a stateless relational query interface of this kind, with lazy exploration of solutions and well defined success, failure and error indications, is not a conformant RELATION node.
%
%\subsubsection{Scenario}
%
%In the scenario depicted in Figure~\ref{fig:browser-to-relation-node}, a user equipped with an ordinary web browser have been talking, via the stateless HTTP API, to the RELATION node at http://n1.org. The web-based terminal displays the remnants of a session during which the user executed two queries. 
%
%\begin{figure}[h]
%    \centering
%	\includegraphics[width=9cm]{browser-to-relation-node-3}
%    \vspace{0.2cm}\caption{A scenario involving a web browser in conversation over HTTP with a remote RELATION node. The two question marks signify that neither the owner-defined resources nor the way answers are produced is known by the user.}
%    \label{fig:browser-to-relation-node}
%\end{figure}
%
%\noindent The interaction started with the loading of the Prolog shell at \texttt{http://n1.org/shell}. This is when the welcome banner as well as the \texttt{?-} prompt inviting the user to a session were displayed. A JavaScript process started running a loop tasked with reading queries from the terminal, transforming them into stateless HTTP requests, picking up responses, and writing answers to the terminal. 
%
%Let us look at how the first solution to \texttt{?-wife(H,W)} was retrieved by using the stateless HTTP API. The URI pointing to this solution is:
%
%\begin{verbatim}
%http://n1.org/call?goal=wife(H,W)&limit=1&offset=0
%\end{verbatim}
%
%\noindent Note the use of \texttt{limit} and \texttt{offset} parameters. Although here they are parameters rather than options, their defaults and the effect that they have on answer terms is the same as for the options accepted by \texttt{toplevel\_call/3}. 
%
%URIs such as these are simple, they are meaningful, they are declarative, they can be bookmarked, and responses are cacheable by intermediates. Most of these desirable properties are there because the client-node interaction involved during the process of resolving the URIs is stateless. 
%
%The form of interaction highlighted in the scenario in Figure~\ref{fig:browser-to-relation-node} can be held over an HTTP connection, making one request for each solution to a given query. When the user entered the query \texttt{?-wife(H,W)}, (and assuming that a relative path was used instead of an absolute one) the following GET request was made to the node:
%
%\begin{verbatim}
%GET /call?goal=wife(H,W)&limit=1&offset=0 HTTP/1.1
%Accept: application/json
%\end{verbatim}
%
%\noindent The value of the \texttt{Accept} header asks the node controller to deliver responses in either JSON or Prolog format. (JSON is the default, so in this case the header could be left out.)\footnote{It might make sense to also support a \texttt{format} parameter as this allows us to easily make a request from the address field of a web browser.}
%
%The node responded with the following JSON structure, indicating that the query succeeded: 
%
%\begin{verbatim}
%{ "type": "success",
%  "data": [{"H":"socrates","W":"xantippa"}], 
%  "more": true }
%\end{verbatim}
%
%\noindent The \texttt{data} property represents the variable bindings resulting from the call, and the value \texttt{true} of the \texttt{more} property signifies that more solutions may be available. This triggered the printing of the variable bindings and served as a prompt for how to continue the exploration.
%
%The user wanted to know if there were more solutions and typed a semicolon which led to the following request being made, looking exactly like the previous request, except that the \texttt{offset} parameter was incremented from \texttt{0} to \texttt{1}:
%
%\begin{verbatim}
%GET /call?goal=wife(H,W)&limit=1&offset=1 HTTP/1.1
%\end{verbatim}
%
%\noindent The node responded like so:
%
%\begin{verbatim}
%{ "type": "success",
%  "data": [{"H":"aristotle","W":"pythias"}], 
%  "more": false }
%\end{verbatim}
%
%\noindent The \texttt{more} property now indicated that no more solutions exist, and after having printed the next set of bindings and a full stop, the shell again produced the \texttt{?-} prompt, inviting the user to enter another query.
%
%At this stage of the interaction, it is reasonable for the user to expect that \texttt{wife/2} has the mode \texttt{wife(?,\,?)} and that queries such as \texttt{?-wife(socrates,\,W)} or \texttt{?-wife(H,\,pythias)} should also succeed. A query such as \texttt{?-wife(plato,\,W)} should fail, however, and the response would be the following:
%
%\begin{verbatim}
%{ "type": "failure" }
%\end{verbatim}
%
%\noindent Finally, an erroneous query, such as the attempt in the scenario to query a predicate which is not defined by the node, should raise an error, and show up as a response like so:
%
%\begin{verbatim}
%{ "type": "error",
%  "data": "Unknown procedure: =/2" }
%\end{verbatim}
%
%\noindent We can summarize this by saying that a RELATION node must offer clients a truly relational query interface over stateless HTTP. Any node failing to meet this expectation is not a conformant RELATION node. A RELATION node is required to respond to a client as if it were Prolog, thereby allowing the user to lazily backtrack over possible solutions. Although a RELATION node may not truly perform backtracking, it must behave externally as though it does. 
%
%\subsection{Browser talking to an ISOBASE node}\label{sec:example-scenario2}
%
%\subsubsection{Profile requirements}
%A conformant ISOBASE node must satisfy all requirements of the RELATION profile and, in addition, make a substantial subset of Web Prolog language features available to clients. In particular, it must offer access, via the stateless HTTP API, to a large set of ISO Prolog predicates, selected standard library predicates, definite clause grammars and predicates for interacting with the Web. Clients may rely on these predicates being present whenever they address a node that advertises the ISOBASE profile.
%
%An ISOBASE node must also support \textit{nondeterministic remote procedure calls} between nodes over stateless HTTP. Concretely, it must provide predicates such as \texttt{rpc/2-3} for synchronous remote calls, and \texttt{promise/3-4} and \texttt{yield/2-3} for asynchronous remote calls whose answers can be retrieved later. These predicates must be implementable on top of the stateless HTTP API of RELATION nodes and must respect the same success, failure and error conventions.
%
%Because all interaction is routed through a stateless request--response protocol, an ISOBASE node must not expose predicates that require persistent interaction state inside an actor, such as the ISO predicates for modifying a private dynamic database (\texttt{assert/1}, \texttt{retract/1} and related predicates) or for interactive input and output (\texttt{read/1}, \texttt{write/1} and similar). When such predicates are called by a remote client, they must behave as if they were not available, typically by raising an \texttt{Unknown procedure} error. An ISOBASE node may offer additional language features or convenience predicates, but clients must assume only the capabilities mandated by the ISOBASE profile and by the underlying RELATION profile.
%
%\subsubsection{Scenario}
%
%In the scenario illustrated in Figure~\ref{fig:browser-to-n1}, the user is running the same query as in the previous scenario, so we do not need to repeat the details of the HTTP requests and responses here:
%
%\begin{figure*}[h]
%    \centering
%	\includegraphics[width=\textwidth]{browser-to-isobase-node-5}
%    \vspace{0.3cm}\caption{Stateless interaction in a scenario involving a web browser and a remote node. }
%    \label{fig:browser-to-n1}
%\end{figure*}
%
%\noindent We focus on the language capabilities of the node and the contents of the shared Prolog database offered by it. The main difference compared to the previous scenario is that an ISOBASE node \textit{must} make a fairly large number of Web Prolog built-in predicates available, including a large subset of ISO Prolog predicates (recall the cloud of ISO Prolog predicates in Chapter~\ref{sec:web-prolog}), standard library predicates, definite clause grammars, as well as predicates for working with the Web. (The scenario differs from the previous one in that it allows the user to browse the source code defining the predicates offered by the owner of the node.)
%
%In this scenario, we make use of \texttt{rpc/2}, an easy-to-use predicate for making \textit{nondeterministic} RPCs from one node to another over stateless HTTP. Note that the URI \texttt{http://n1.org} in the first argument of the call to \texttt{rpc/2} shows that it accesses the relation \texttt{wife/2} made available by the node \texttt{n1.org} in the previous scenario. Of course, the user could have made the query like so instead:
%
%\begin{verbatim}
%?- rpc('http://n1.org', wife(H, W)).
%H = socrates, W = xantippa ;
%H = aristotle, W = pythias.
%?-
%\end{verbatim}
%
%\noindent Asynchronous RPC via \texttt{promise/3-4} and \texttt{yield/2-3} is also available. It turns out that these RPC predicates may all be implemented on top of the stateless HTTP API. They will be dealt with in more detail in Chapter~\ref{sec:the-prolog-web}.
%
%A stateless protocol comes with certain limitations. One thing we cannot do is that once we have entered a query, there is no way we can abort it. In fact, no outside party, including the client who originally submitted the request, can communicate in any way with the remote computation once it has begun evaluation. A Web Prolog query may take a long time to execute, or may never return with an answer, and yet, once a query has been entered, it cannot be aborted by the user. Instead, queries must always be executed under a time limit so that, typically after just a few seconds or even less, a client would be met with an error message. This is what happened with the user's second query: it encountered a \texttt{Timeout\,exceeded} message. Browsing the contents of the node's shared database, the problem seems to be that the node owner forgot to add a DCG rule \texttt{s\,{-}->\,[]} that would have allowed the call to terminate.
%
%Another limitation is that a truly stateless request--response protocol simply cannot support updates to the dynamic private database of an actor, nor can it support I/O. It means that ISO predicates for modifying the dynamic database (e.g. \texttt{assert/1} and \texttt{retract/1}), or I/O (e.g. \texttt{read/1} and \texttt{write/1}), are \textit{not} available. %As we shall see, it is the nature of the protocol that constrains what \textit{can} be made available. 
%This is what happened with the user's third query, which was met with an \texttt{Unknown\,procedure} error message when trying to assert the clause \texttt{p(a)}. The predicate \texttt{assert/1} is simply not available to a client talking to an ISOBASE node.
%
%\subsection{Browser talking to an ISOTOPE node}\label{sec:browser-talking-to-isotop-node}
%
%\subsubsection{Profile requirements}
%
%A conformant ISOTOPE node must satisfy all requirements of the ISOBASE profile and, in addition, offer clients access to long-lived toplevel actors via a semi-stateful HTTP API (without directly exposing the internal \texttt{toplevel\_*} predicates). It must provide operations for spawning remote toplevel actors (for example via a \texttt{/toplevel\_spawn} endpoint) and for sending queries to, and retrieving answers from, an existing toplevel (for example via \texttt{/toplevel\_call} and \texttt{/toplevel\_next}). Each such interaction remains an ordinary stateless HTTP request, but the requests must carry a session key (a process identifier, pid) that identifies the toplevel actor whose private state they address.
%
%Within a session associated with a toplevel actor, an ISOTOPE node must allow clients to treat that actor as a conventional Prolog process with a private dynamic database and support for interactive input and output. In particular, predicates for updating the dynamic database, such as \texttt{assert/1} and \texttt{retract/1}, must be available and must affect only the private workspace of the toplevel actor identified by the pid. Likewise, certain I/O predicates may be redefined in terms of primitive operations such as \texttt{output/1} so that they can deliver output to the client as part of the HTTP interaction sequence.
%
%An ISOTOPE node must also provide a way for clients to abort long running or non terminating queries addressed to a toplevel actor and to terminate inactive sessions, for example by interpreting control signals (such as Control C-in a browser shell) as explicit abort requests and by applying idle timeouts to actors that have seen no activity for some configured period. The precise policies for timeouts and resource reclamation are left to the node implementation, but a conformant ISOTOPE node must ensure that long-lived toplevel sessions remain manageable and that clients can reliably start, use and eventually discard such sessions.
%
%\subsubsection{Scenario}
%
%In the scenario in Figure~\ref{fig:browser-to-isotope-node}, the client has established communication with a toplevel actor running on an ISOTOPE node, which provides all the capabilities of an ISOBASE node, and then some.
%
%\begin{figure}[h]
%    \centering
%	\includegraphics[width=\textwidth]{browser-to-isotope-node-3}
%    \vspace{0.1cm}\caption{A scenario involving a web browser in conversation over semi-stateful HTTP with a remote ISOTOPE node running a toplevel actor.}
%    \label{fig:browser-to-isotope-node}
%\end{figure}
%
%\noindent The request below instructed the node controller to spawn a toplevel actor running as a session, and to populate its private database with the value of the \texttt{load\_text} option. (Let us assume that the clause \texttt{q(X)\,:-\,p(X)} was picked up from an editor window that is not shown in the diagram.)
%
%\begin{verbatim}
%POST /toplevel_spawn HTTP/1.1
%Content-Type: application/json
%
%{ "options": "[session(true), load_text('q(X):-p(X).')]" }
%\end{verbatim}
%
%\noindent The node responded with a JSON structure containing the pid for the spawned toplevel actor:
%
%\begin{verbatim}
%{ "type": "spawned",
%  "pid" : 56120932 }
%\end{verbatim}
%
%\noindent The pid must be sent along when querying the toplevel:
%
%\begin{verbatim}
%GET /toplevel_call?pid=56120932&goal=q(X)&limit=1&offset=0 
%\end{verbatim}
%
%\noindent Note that except for the pid, the JSON structure representing this solution is identical to the response when making a stateless request:
%
%\begin{verbatim}
%{ "type": "success",
%  "pid" : 56120932,
%  "data": [{"X":"a"}],
%  "more": true }
%\end{verbatim}
%
%\noindent When the user asked for another solution, the following request was made:
%
%\begin{verbatim}
%GET /toplevel_next?pid=56120932 HTTP/1.1
%\end{verbatim}
%
%\noindent Note that it took the \textit{combination} of the clause that was loaded into the toplevel's private Prolog database, and data in the shared database to answer the query \texttt{?-q(X)}.
%
%It makes sense to allow the actor's dynamic database to be updated by means of \texttt{assert/1} and \texttt{retract/1}:
%
%\begin{verbatim}
%GET /toplevel_call?pid=56120932&goal=assert(r(c,d)) HTTP/1.1
%\end{verbatim}
%
%\noindent With this, the clause \texttt{r(c,d)} ended up in the private database of \texttt{56120932}, but was removed again when the user decided to call \texttt{retract(r(c,d))}.
%
%The \texttt{writeln/1} predicate has been redefined by the node controller as follows:
%
%\begin{verbatim}
%writeln(Term) :- output(Term).
%\end{verbatim}
%
%\noindent The user calling \texttt{greetings/0} resulted in a series of messages of the following form to be sent back to the client, one per second:
%
%\begin{verbatim}
%{ "type": "output",
%  "pid" : 56120932,
%  "data": "hello" }
%\end{verbatim}
%
%\noindent For this to work over HTTP, a trick somewhat similar to so-called \textit{long-polling} may be employed: every time a JSON structure of type \texttt{output} reaches the client, it must issue a new request for a response, which can come as more output, or some other response.
%
%Note that the call to \texttt{greetings/0} will run forever unless the user does something about it. When the user hits Control-C, the nonterminating query is aborted, and the node responds with the following message:
%
%\begin{verbatim}
%{ "type": "abort",
%  "pid" : 56120932 }
%\end{verbatim}
%
%\noindent If the client does not exit the remote process, the node may terminate it after a set time with no activity. The pid will no longer be valid and any changes made to the workspace of the process will be lost. At least this is one way of handling it.
%
%This creates an opportunity for a node implementer to support persistency of a simple kind. Depending on a node setting, just before terminating, all code residing in the workspace of the process can be saved to a file (i.e. ``marshalled'') or to another form of persistent storage (e.g. under a name identical to the pid) before the process is finally terminated.
%
%Should a query request targeting this pid be made later (possibly much later), the node can spawn the process again (possibly give it the same pid again) and inject it with the content of the stored file. This would allow a client to hang on to a particular toplevel over an extended period of time and still not demand too many resources from the node.
%
%\subsection{Browser talking to an ACTOR node}\label{sec:browser-talking-to-node}
%
%\subsubsection{Profile requirements}
%A conformant ACTOR node must satisfy all requirements of the ISOTOPE profile and, in addition, expose the full set of Erlang-style concurrency primitives of Web Prolog to clients. In particular, it must allow clients to spawn new actors (for example via \texttt{spawn/2} or related predicates), send and receive messages between actors (for example via \texttt{!/2} and \texttt{receive/1}), and use additional primitives such as \texttt{self/1}, monitoring and linking operations and termination operations (such as \texttt{exit/2}) to structure supervision and failure handling. These predicates must support concurrent execution of many actors within the node and must make actor communication appear network transparent where appropriate.
%
%On the interaction side, an ACTOR node must provide a stateful, bidirectional WebSocket API in addition to the HTTP APIs of the underlying profiles. Over this WebSocket connection, the node must be able both to receive commands from the client (such as requests to spawn toplevel actors or ordinary actors, to submit queries, or to send messages) and to push messages and results back to the client without requiring a new HTTP request each time. This enables genuinely mixed initiative interaction patterns in which both the client and the node may take the initiative at different times.
%
%An ACTOR node may still provide HTTP based access to certain operations, including the semi-stateful toplevel API of ISOTOPE, but clients may rely on the presence of a WebSocket API for rich, low latency interaction with actors running inside the node. As with the other profiles, an ACTOR node may implement more capabilities than the profile mandates, but any node that does not offer clients both the language level actor primitives and a corresponding stateful WebSocket based interaction channel is not a conformant ACTOR node.
%
%\subsubsection{Scenario}
%
%Figure~\ref{fig:browser-to-toplevel} illustrates a scenario in which a user engages with an ACTOR node over a stateful WebSocket connection.
%
%\begin{figure}[h]
%    \centering
%	\includegraphics[width=\textwidth]{browser-to-pengine-4}
%    \caption{A scenario involving a web browser in conversation with a remote node running a toplevel actor over a stateful WebSocket connection.}
%    \label{fig:browser-to-toplevel}
%\end{figure}
%
%\noindent A description of the step-by-step evolution of the scenario in Figure~\ref{fig:browser-to-toplevel} can be given as follows. After having loaded the shell and established a WebSocket connection with the node, the JavaScript process sent the following command to the node's controller, instructing it to spawn a toplevel actor running a session with the clause \texttt{q(X)\,:-\,p(X)} loaded into the actor's private database:
%
%\begin{verbatim}
%{ "command": "toplevel_spawn",
%  "options": "[session(true), load_text('q(X) :- p(X).')]" }
%\end{verbatim}
%
%\noindent As requested by the client, the pid of the toplevel actor was returned to the client, not in the form of an HTTP response, but embedded in a JSON formatted message sent over the WebSocket connection:
%
%\begin{verbatim}
%{ "type": "spawned",
%  "pid" : 23981144 }
%\end{verbatim}
%
%\noindent This produced the prompt \texttt{?-} in the terminal, inviting the user to enter a query. The query \texttt{?-q(X)} caused the JavaScript shell to send the following message to the node over the WebSocket connection:
%
%\begin{verbatim}
%{ "command": "toplevel_call",
%  "pid": 23981144,
%  "goal": "q(X)",
%  "options": "[limit(1)]" }
%\end{verbatim}
%
%\noindent The node controller converted it to a \texttt{toplevel\_call/3} call and converted the answer term into a JSON structure encoding the first solution: \todo{Perhaps this should move to the ISOTOPE scenario?}
%
%\begin{verbatim}
%{ "type": "success",
%  "pid" : 23981144,
%  "data": [{"X":"a"}],
%  "more": true }
%\end{verbatim}
%
%\noindent The value \texttt{true} of the \texttt{more} property signifies that more solutions may be available. Asking for the next solution results in a similar message, except that the value of the \texttt{data} property becomes \texttt{[\{"X":"b"\}]} and the value of \texttt{more} becomes \texttt{false}.
%
%Knowing that the node's shared database offers a definition of \texttt{count\_actor/1}, the user then decided to spawn a count actor of the kind we saw in Chapter~\ref{sec:web-prolog}. Here is the message sent over the WebSocket connection:
%
%\begin{verbatim}
%{ "command": "toplevel_call",
%  "pid": 23981144,
%  "goal": "spawn(count_actor(0),Pid),register(counter,Pid)",
%  "options": "[limit(1)]" }
%\end{verbatim}
%
%\noindent This resulted in:
%
%\begin{verbatim}
%{ "type": "success",
%  "pid" : 23981144,
%  "data": [{"Pid":91883433}],
%  "more": false }
%\end{verbatim}
%
%\noindent While the toplevel actor continues to serve shell queries, the spawned count actor runs concurrently and receives \texttt{count} messages independently of the shell interaction, illustrating the concurrent, actor based nature of the ACTOR profile. With access to the pid of the spawned count server, the user could send it a \texttt{count} message in order to try it out, and call \texttt{flush/0} to inspect the response.
%
%Finally, the user called \texttt{toplevel\_spawn/3} with the \texttt{node} option. This is yet another way to do network-transparent concurrency that will be dealt with in more detail in Chapter~\ref{sec:the-prolog-web}.
%
%Should the user leave the terminal at \texttt{http://n7.org/shell} the toplevel process will first kill (by calling \texttt{exit/2}) the count actor and then die (since the closing of the WebSocket will kill it, again by a call to \texttt{exit/2}). There is a kind of supervision hierarchy in play here, and we shall say more about that in Chapter~\ref{sec:the-prolog-web}.
%
%\subsection{Social robot talking to an ACTOR node}\label{sec:robot-talking-to-actor-node}
%
%In the previous sections, the toplevel actor at the core of the shell is the root of control -- the first process or supervisor from which others are spawned or monitored. However, the root of an application is not always a toplevel actor - it may be any kind of actor.
%
%--
%
%Here is another scenario: As suggested in Figure~\ref{fig:robot-to-actor-node}, we can imagine a social robot interacting with the echo server over the stateful web API. It just echoes the input back to the person interacting with it (perhaps using speech in both directions). For this to work, the robot needs to be able to create an actor on a remote node and access it over an API using WebSockets as transport. 
%
%\begin{figure}[h]
%    \centering
%	\includegraphics[width=11.5cm]{robot-to-actor-node}
%    \caption{Stateful interaction in a scenario involving a social robot and a remote node. }
%    \label{fig:robot-to-actor-node}
%\end{figure}
%
%\noindent Although it makes no difference in this particular case, statefulness in general allows the build-up of a context which makes earlier interactions matter to how later ones are handled. For this to work, the node needs to maintain one such context per client session and therefore must be able to distinguish one client from another. As we have seen, the pid acts as an identifier that lets a node on the Prolog Web do just that.
%
%\noindent Facial expressions are important in the communication process as they help reinforce a positive or negative view of a message, help the sender reinforce the message being sent, and help the receiver understand the message being received.
%
%---
%
%Importance of timing and real-time asynchronous interaction and concurrency.
%
%---
%
%In this scenario, the actor spawned by the robot client is not a toplevel actor. Rather, it is an actor implementing a so called \textit{dialog manager} (or \textit{dialog controller}), a component in a conversational system that is responsible for ... The design of such a component differ from systems to systems. 
%
%\begin{verbatim}
%{ "command": "spawn", "goal": "dialog_manager" }
%\end{verbatim}
%
%---
%
%event-driven statemachine

%SCXML
%
%
%\section{SCRAP: Node profiles}\label{}
%
%\todo{Some of this might be useful, but must me compressed.}
%
%In the context of W3C standards, a \textit{profile} generally refers to a subset of the full standard that is tailored to meet specific requirements or constraints, often relevant to a particular domain or application. Profiles are crafted to simplify implementation or to optimize performance, sometimes by omitting certain features that are deemed unnecessary for the specialized use-case at hand. 
%
%The language and interactional capabilities of a node profile determines a Web Prolog language fragment as well as one or more communication protocols and web APIs that support them. The major rationales for introducing profiles is that the less capable nodes will be easier to standardize and easier to implement than the more capable ones, but there are also other (minor) reasons why having profiles is a good idea.
%
%The Venn diagram in Figure~\ref{fig:profiles} gives a overview of the profiles supported by nodes in the proposed design of the Prolog Trinity ecosystem. Four profiles will be presented and discussed in some detail: RELATION, ISOBASE, ISOTOPE and ACTOR. They form a hierarchy in the sense that nodes implementing the ACTOR profile are more capable than nodes implementing the ISOTOPE profile, and nodes implementing the ISOTOPE profile are more capable than nodes implementing the ISOBASE profile. The RELATION profile is the least capable of these four.
%
%For example, the reason that owner-defined predicates and the HTTP API are located inside the area marked RELATION is that it signifies that those are features supported by \emph{all} profiles, and the fact that the WebSocket API is placed in the outermost area implies that it is only available in the ACTOR profile.
%
%\begin{figure*}[tb]
%    \centering
%	\includegraphics[width=10cm]{profile-hierarchy-new}
%    \caption{Node profiles.}
%    \label{fig:profiles}
%\end{figure*}
%
%In short, running against a node that offers the ACTOR or ISOTOPE profile, we can both query and (through source code injection) program a toplevel, but the ACTOR profile alone allows us to use predicates such as spawn, send, receive, assert and retract in queries or programs. Running against an ISOBASE or RELATION node, we can only query it, and only the ISOBASE profile allows us to pose queries using built-in predicates. In this book, by far the most space has been devoted to the presentation of the ACTOR profile. 
%
%---
%
%The suggested hierarchy of profiles should be regarded as tentative, and four profiles may not be enough to cover all needs, so towards the end of the chapter we shall also sketch a couple of (less worked out) ideas for other profiles.
%
%\subsection{Capabilities}\label{}
%
%As should be fairly clear from Figure~\ref{fig:profiles}, the language capabilities of the different profiles varies a lot. A node with the RELATION profile does not offer any built-in predicates at all, the ISOBASE and ISOTOPE nodes come with a large set of ISO Prolog built-in predicates, and the ACTOR profile with an even larger one.
%
%As for predicates useful for distributed programming, the ISOBASE and ISOTOPE profiles comes with only one, namely \texttt{rpc/2-3}. 
%
%---
%
%M. V. Hermenegildo et al.: "This modular design allows moving from pure LP, where, e.g., no impure built-ins are visible, to full ISO Prolog by specifying the set of imported modules, and going well beyond that into a multi-paradigm language, while maintaining full backwards compatibility with standard Prolog."
%
%\subsection{The ACTOR profile}\label{}
%
%The capabilities of the ACTOR profile have already been described and thoroughly demonstrated in previous chapters, and this chapter describes the other and less capable profiles mainly in terms of what we can no longer expect from nodes that offer them. A brief summary of the capabilities of an ACTOR node is in order, though.
%
%As it supports both sequential distributed programming over stateless HTTP as well as concurrent distributed programming over stateful WebSockets, the ACTOR profile is the most powerful Web Prolog profile. The column to the left in the area corresponding to the ACTOR profile lists some of the predicates that are inspired by Erlang. The column in the middle lists predicates for creating and communicating with Prolog toplevels. The third column is interesting mostly because it shows what kind of predicates -- for I/O and dynamic database manipulation -- are \textit{not} available in the other three profiles. 
%
%A conforming ACTOR node is definitely the most difficult to implement. 
%
%\subsection{The ISOTOPE profile}\label{sec:isotope-profile}
%
%In terms of interactional as well as language capabilities, the ``distance'' from the ACTOR profile to the ISOTOPE profile is quite big. An ISOTOPE node does not support the WebSocket protocol and therefore none of the interactional capabilities relying on stateful and asynchronous inter-node communication can be offered by such a node. Only the stateless and synchronous communication offered by the stateful HTTP API is available, and an ISOTOPE node therefore does not allow a client to create a session on the node -- only one-shot queries are possible. Consequently, a node that offers the ISOTOPE profile lacks the language capabilities that has to do with spawning and messaging. Predicates such as \texttt{spawn/1-3}, \texttt{!/2}, \texttt{receive/1-2} and the \texttt{toplevel\_*} predicates are simply not defined by a node with this profile. An attempt to call any of these predicates, in an injected program or in a query, will result in a runtime error. 
%
%Regarding the \texttt{rpc/2-3} predicate, the \texttt{limit} option as well as the \texttt{load\_*} options are supported by a node with the ISOTOPE profile. Indeed, as we shall see, when calling \texttt{rpc/2-3} the \texttt{limit} option is always available, regardless of profile, but the \texttt{load\_*} options are not available in profiles less capable than the ISOTOPE profile.
%
%Recall that \texttt{rpc/2-3} implements a form of synchronous RPC. As it uses HTTP as transport, asynchronous RPC using \texttt{promise/3-4} and \texttt{yield/1-2} can also be made to an ISOTOPE node. %The reason for imposing this restriction is that we want to ensure the less capable profiles can be implemented in languages not supporting concurrency, and that an implemention of \texttt{promise/3-4} and \texttt{yield/1-2} seem to require some kind of concurrency at the client side.\footnote{This is not strictly true. JavaScript is not a concurrent language but can still support so called \textit{futures}, which is just another name for promises.}
%
%The only way to inject source code into a remote process is to use the \texttt{load\_*} options provided by \texttt{rpc/2-3} and \texttt{promise/3-4}. Predicates for manipulating the dynamic database -- \texttt{assert/1} and \texttt{retract/1} -- are not available either. Using them would make little sense anyway, since a session cannot be created.
%
%\subsection{The ISOBASE profile}\label{sec:isobase-profile}
%
%\noindent %The ISOBASE profile requires that a conforming node can be queried using a subset of ISO Prolog. We refrain from trying to determine the size of this subset at this point, but as we have already mentioned, it does not include predicates for I/O, database manipulation or user-defined operators. An ISOBASE node also offers querying over stateless HTTP, except that the \texttt{load\_*} parameters are not supported. The \texttt{rpc/2-3} predicate is supported, but only over the HTTP transport. and nor there is source code injection possible.
%In terms of capabilities, the distance from ISOTOPE to ISOBASE is fairly small. Only a few built-in predicates present in the ISOTOPE profile are missing from the ISOBASE profile.
%
%For a non-Prolog client only the HTTP API is available here too, and its powers are somewhat reduced since the \texttt{load\_text} query parameters for code injection is not supported. Because of this lack of support for source code injection from the underlying HTTP transport, \texttt{rpc/2-3} does not support the \texttt{load\_*} options either:
%
%\begin{verbatim}
%?- rpc('http://isobase.org', foo(X), [
%       load_text('foo(a). foo(b).')
%   ]).
%Error: The ISOBASE profile does not support this option
%?-
%\end{verbatim}
%
%\noindent Moreover, an ISOBASE node is not required to support \texttt{op/3}, the ISO Prolog predicate allowing a programmer to define prefix, postfix or infix operators and specify their precedences. After all, \texttt{op/3} is normally used as a directive, so to offer it in a profile that does not also support the injection of source code makes little sense. More importantly, user-defined operators makes writing a parser for Web Prolog more difficult, and the most important motivation for having profiles in the first place is to make it easy to implement the less capable profiles.
%
%\subsection{The RELATION profile}\label{sec:relation-profile}
%
%In terms of language capabilities, the distance from the ISOBASE profile to the RELATION profile is again considerable. A node supporting the RELATION profile is not required to define any built-in predicates \textit{at all}. Not even control predicates such as \texttt{,/2} or \texttt{;/2} are available. Only predicates explicitly exported by a node-resident program can be queried. If they are not exported, they will be treated as undefined. 
%
%In terms of interactional capabilities, a RELATION node offers a non-Prolog client the same stateless HTTP API as an ISOBASE node, complete with options such as \texttt{offset} and \texttt{limit}. Consequently, a Web Prolog client can query a RELATION node using \texttt{rpc/2-3}.
%
%However, because a RELATION node does not offer any built-in Web Prolog predicates, the following will not work even if \texttt{p/1} and \texttt{q/1} are defined by the node-resident program (and unless \texttt{;/2} is defined by the owner):
%
%\begin{verbatim}
%?- rpc('http://relation.org', (p(X);q(X))).
%ERROR: Undefined procedure: ;/2
%?-
%\end{verbatim}
%
%\noindent The main motivation for allowing a node to offer clients as little as this is to make it relatively easy to implement a Prolog Web node in just about any programming language. At the very least, an implementation must be able to parse Prolog terms, as well as to be able to create bindings for the free variables in a query. %It does not necessarily need to do any backtracking or recursion but it must \textit{behave} as if it does. 
%
%A RELATION node that offers an HTTP API with a fixed set of owner-defined predicates may be compared to a traditional web API and to variants of remote procedure calls on the Web such as XML-RPC, JSON-RPC and SOAP that only offer what amounts to a particular call to a node-resident program. At first blush, they seem to be equally inflexible, but only until we realise that when querying a RELATION node we are dealing with a \textit{relational} language.
%
%When querying predicates with an arity greater than 1, we can submit queries such as \texttt{?-r(a,b)}, \texttt{?-r(a,Y)}, \texttt{?-r(X,b)} and \texttt{?-r(X,Y)}. In other words, provided the mode of \texttt{r/2} is unrestricted, the API allows us to ask different \textit{kinds} of questions. This is the power of a relational query interface, a power a more traditional HTTP API cannot provide clients with. Likely, a traditional HTTP API would force us to introduce different endpoints for different kinds of queries, something that makes an API more complicated.
%
%Unless a RELATION node offers a predicate that does not (always) terminate in finite time, it also safeguards against nonterminating queries, a protection nodes offering the other profiles do not give. Furthermore, making sure all predicates that are offered by the node are deterministic ensures that no application state has to be maintained by the node and thereby eliminates the need for a cache mechanism such as the one described in Section~\ref{sec:efficient-implementation-of-the-http-api}. 
%
%---
%
%\noindent With advanced features such as Erlang-style concurrent programming with actors, the most capable Web Prolog profile is not trivial to implement. This threatens to block some of the current Prolog systems creating the wrapper necessary for the realization of ...
%
%To this end, we propose a \textit{profile-based} approach, where different Prolog systems can offer different levels of capabilities on the Prolog Web, depending on their strengths. We envision a \textit{hierarchy} of profiles, each corresponding to a particular level of functionality a system can provide, from basic ISO Prolog compatibility to more advanced features like tabling, constraint solving, or concurrency.
%
%For instance, a basic profile may provide very little beyond the standard ISO Prolog core, whereas more advanced profiles may support tabling, constraints, or a sophisticated concurrency model. Each Prolog system can then implement the profile that best matches its capabilities, allowing it to participate in the Prolog Web without the need to match every feature that other systems may provide.
%
%This hierarchical approach also facilitates gradual improvements. As more Prolog systems implement higher-level profiles, they expand their ability to interact and contribute more complex services to the Prolog Web. In this way, a Prolog system like Scryer Prolog, which is lightweight and ideal for experimentation, can still interact meaningfully with systems like SWI-Prolog, which offer an extensive standard library and interfaces to external systems.
%
%The idea of profiles ensures that even the smallest, simplest Prolog implementations can establish their presence on the Prolog Web, while more feature-rich systems can offer more complex services. This approach enables cooperation across the Prolog spectrum, establishing a consistent foundation for future innovation.
%
%---
%
%A profile of Prolog that caters for both portability and interoperability.
%
%Actually, we shall present a \textit{hierarchy} of profiles as well as rules for how programs written in one profile may interact with a program written in another profile. (After all, interoperability is not necessarily symmetric.)
%
%After all, most of the systems in Figure~\ref{fig:prolog-systems-logos} adhere to the ISO Prolog core standard. With the addition of a standard library, the community might be able to come up with a standardized language which allows programs written in it to be ported between systems.
%
%The ``dumbing down'' of a language, if done in the right way and for the right reason, can actually be smart.
%
%---
%
%The pid of the created toplevel is returned to the client, not in the form of the binding of a variable (as in the predicate API), but as a message of the form \texttt{\{"type":"spawned","pid":<pid>\}}. 
%
%The full set of options valid in the predicate API is valid in the web API as well. Since the options for \texttt{toplevel\_spawn/2},  \texttt{toplevel\_call/3} and \texttt{toplevel\_next/2} are documented in Appendix~\ref{sec:toplevel-api}, we do not give the details here. A brief reminder should be sufficient. In the message corresponding to \texttt{toplevel\_spawn/2} we can use options such as \texttt{exit}, \texttt{monitor}, \texttt{link}, \texttt{timeout}, \texttt{load\_list}, \texttt{load\_text}, \texttt{load\_uri} and \texttt{load\_predicates}. (Recall that the majority of these options are inherited from \texttt{spawn/3}.) For \texttt{toplevel\_call/3} we can use \texttt{template}, \texttt{offset} and \texttt{limit}. For \texttt{toplevel\_next/2} there is only one valid option, namely \texttt{limit}.
%
%---
%
%\noindent %The ISOBASE profile requires that a conforming node can be queried using a subset of ISO Prolog. We refrain from trying to determine the size of this subset at this point, but as we have already mentioned, it does not include predicates for I/O, database manipulation or user-defined operators. An ISOBASE node also offers querying over stateless HTTP, except that the \texttt{load\_*} parameters are not supported. The \texttt{rpc/2-3} predicate is supported, but only over the HTTP transport. and nor there is source code injection possible.
%In terms of capabilities, the distance from ISOTOPE to ISOBASE is fairly small. Only a few built-in predicates present in the ISOTOPE profile are missing from the ISOBASE profile.
%
%For a non-Prolog client only the HTTP API is available here too, and its powers are somewhat reduced since the \texttt{load\_text} query parameters for code injection is not supported. Because of this lack of support for source code injection from the underlying HTTP transport, \texttt{rpc/2-3} does not support the \texttt{load\_*} options either:
%
%\begin{verbatim}
%?- rpc('http://isobase.org', foo(X), [
%       load_text('foo(a). foo(b).')
%   ]).
%Error: The ISOBASE profile does not support this option
%?-
%\end{verbatim}
%
%---
%
%In cases where a node hosts programs and data, external users are not able to make any changes to them, only the owner of the node is allowed to do so. However, as was demonstrated in Section~\ref{sec:private-database}, client programmers are allowed to complement them, and occasionally shadow them, with predicates that they themselves have written.
%
%It follows that in cases where a node does \textit{not} host such programs and data, it can still serve as a \textit{compute server}, allowing authorized clients to supply it with user-defined predicates, followed by queries to be executed over them. (This provides a suitable foundation for the so called \textit{Prolog playgrounds} that will be discussed in Chapter~\ref{sec:ides}.)%[TODO: Move this to text about source code injection?]
%
%If no such predicates are loaded by the client, the answers to subsequent queries will depend only on predicates defined and loaded by the owner in combination with the built-in Web Prolog predicates supported by the node. A relevant use case might be a database application with knowledge we would like to share, knowledge that may be backed up by a relational database or a semantic web triple store.
%
%When a node hosts some resident predicates, and a client injects more predicates, the resulting program is formed by the union of the resident and the injected predicates. However, should a predicate in the private database be present also in the node-resident shared database, the definition of the former will shadow the latter.
